\section{Quadratic inward and outward section estimates}
\label{app:moving-section-bookkeeping}

We estimate the actual quadratic passages on their original sections,
keeping the coefficient-dependent first multipliers. The inward maps
provide heights on one common centre section; their equality defines the
match. The outward maps carry the two lips to the original raw ports.
Together these estimates control the section terms in physical curvature.
The final paired estimate bounds the complete derivative above a
comparison orbit, where the separate first variations may be large.
We use the scalar variation and continuation formulas of the physical
section toolkit throughout.

\subsection{Quadratic radial embeddings on common raw sections}

For the quadratic passage estimates below, fix the original raw-section and
ordinary \(C^3\) observer reserves and the passive compact set.
Phase derivatives freeze \(g,k,\pi,A,C,m\), with \(g,k>0\) independently
small. In the positive or reflected negative radial coordinate, use
\begin{equation}\label{deg:quadratic-inward-family}
 \begin{aligned}
 P_s(v)&=1-s\pi v+3v^2+skv^3,& B_s(v)&=1+g(2+skv),\\
 Q_s(v)&=1-Ag(1+skv)+Cg^2(1+skv)^2,&
 \zeta_s(v)&=sCkg^3(1+skv),\qquad s=\pm1.
 \end{aligned}
\end{equation}
The exact scalar graph is
\begin{equation}\label{deg:quadratic-inward-scalar-common}
 Y_v=\frac{Y(vB_s+\zeta_sY)}{g(P_s+Q_sY)},
 \qquad x=sgkv,\quad z=g^3k^2Y.
\end{equation}
Decreasing \(v\) follows positive time for \(s=+1\) and its inverse
for \(s=-1\). On each domain where the actual lip continuation
exists, denote its original-centre height by \(Z_s(x_s)\) and define
\begin{equation}\label{deg:quadratic-inward-memory-common}
 \begin{gathered}
 A_s(x_s)=\int_0^{x_s}\frac{vB_s}{P_s}\,dv,\quad
 \mathcal D_s=\frac{P_s(x_s)}{x_sB_s(x_s)}\partial_{x_s},\\
 \log Z_s=\log\{\ell_0P_s(x_s)\}-A_s(x_s)/g+\mathcal R_s,
 \quad H_s=P_s'/(vB_s),\quad N_s=-g\mathcal D_s\log Z_s .
 \end{gathered}
\end{equation}
The polynomial factor and action in this formula retain the entrance
height and the full linear contraction to the centre.
Coefficient-only memory remains in \(\mathcal R_s\); the estimates
below concern its positive-order rows. Write \(J=A_+\) and omit
\(s=+1\) subscripts when only the positive graph is used.

On a fixed compact segment \([0,b]\) with \(P_s\ge p_b>0\),
the scalar graph has the common preparation
\begin{equation}\label{deg:positive-large-S-inward-compact-graph}
 \begin{gathered}
 Y_v=\frac{vB_s}{gP_s}Y+K_s(v,Y)Y^2,\qquad
 K_s=\frac{\zeta_sP_s-vB_sQ_s}{gP_s(P_s+Q_sY)},\\
 m_s(v)=\exp\left[-\frac1g\int_v^b\frac{wB_s(w)}{P_s(w)}\,dw\right]
 \le1,\qquad m_s(b)=1 .
 \end{gathered}
\end{equation}
Fix a prescribed height buffer \(0\le Y\le h_b\) before shrinking
\(g,k\). The positive units \(B_s,Q_s\), the margin \(p_b\), and
rational differentiation give
\( |\partial_Y^jK_s|\le C_b/g\), \(0\le j\le3\), on that buffer.
Lemma~\ref{deg:lem-compact-height-continuation} therefore applies with
\(\kappa_c=g\) whenever the incoming height satisfies
\(0<y\le h_b/2\) and \(y/g\) is small. Each regional proof supplies
these input bounds and the logarithmic phase rows of \(y\); the
continuation keeps the same original centre and uses no reciprocal
bound on \(m_s(0)\).

\begin{lemma}\label{deg:lem-quadratic-outer-jets}
Fix the raw section buffers, then a compact neighborhood of the passive
coefficients \(A,C\), and \(\pi\in[\pi_\flat,\Pi]\).
There are fixed \(M,L,d_q,\delta_q,g_0,k_0>0\) such that the following
holds for independently small \(0<g<g_0\), \(0<k<k_0\).
For either \(s=\pm1\), let
\[
 x_0\ge2M,\quad q=x_0^{-1},\quad S=kx_0\le\delta_q,
 \quad a=gkx_0,\quad b_0=-x_0^2/2 .
\]
The orbit from the lip at \(v=g+sa\) reaches \(v=g+sd_q\)
with positive height. Its raw embedding \(H_s(b_0)\), and any increasing
observer \(\Phi(H_s(b_0))\) on the prescribed buffered raw image for which
\(\|\partial_z\log\Phi'\|_{C^1}\le C\), have positive first derivative. Put
\[
 E=x_0\partial_{x_0}=a\partial_a,\qquad
 h_s=\frac{3-s\pi q+q^2+sS}{1+2g+sgS},\quad D_s=h_sE,
 \quad\lambda=\frac{12g}{1+2g},\quad L_g=1+|\log g|.
\]
After separate affine normalizations of the embeddings, their slopes obey
\begin{equation}\label{deg:quadratic-outer-weighted-jets}
 D_s\log\phi_s'=-\lambda h_s+G_s,
 \qquad |G_s|+|D_sG_s|\le CgL_g(q+S).
\end{equation}
Primes use the lip label \(b_0\); \(g,k,\pi,A,C\) are held fixed.
The two raw passages exist separately. The conclusion does not require
their inward central images to intersect.
\end{lemma}

\begin{proof}
The common spectral estimate separates the exact homogeneous power
from its logarithmic residual. We first check outward existence and
admission to the raw buffers, then apply its derivative bounds with
the original moving initial data and terminal observer.

Write \(x=|v-g|\), \(W=z/x^2\), \(\tau=\log(x/a)\), and set
\begin{align*}
 m_s(x)&=3g+sx-s\pi kg^2/x+k^2g^3/x^2,\\
 Q_s(x)&=1-A(g+sx)+C(g+sx)^2,\\
 N_s(x,W)&=1+2g+sx+sCx(g+sx)W,
 \qquad D=m_s+Q_sW .
\end{align*}
Division of the exact raw field gives
\begin{equation}\label{deg:quadratic-outer-ratio}
 W_\tau=F_s(x,W):=W(N_s/D-2),\qquad
 W(0)=\ell_0gp,\quad p=3-s\pi q+q^2+sS.
\end{equation}
Use the exact signed dictionary for Lemma~\ref{deg:lem-spectral-exterior},
\[
 t=s,\quad b=2g,\quad v=g,\quad
 (\mu_2,\mu_1,\mu_0)=(3g,-s\pi kg^2,k^2g^3),\quad
 a=gS=gkx_0,\quad\sigma=gk.
\]
Choose \(M\) large and \(\delta_q\) small. For \(S\le\delta_q\),
\(m_s(x)\asymp g\) on \([a,2a]\), for both signs, and
\[
 |b|+|v|+|\mu_2|+a+|\mu_1|/a+|\mu_0|/a^2\le Cg,
 \qquad \Theta=g(S+\pi q+q^2)\le Cg(S+q).
\]
Thus the common lemma's intrinsic construction applies with
\(\rho\asymp g\). For the negative sign it gives the original
root-free outward passage and invariant strip even when \(m_-\)
changes sign beyond \(2a\).

\smallskip
\noindent\emph{Common positive admission.}
For \(s=+1\), the same construction applies on the larger domain
\begin{equation}\label{deg:quadratic-positive-admission}
 x_0\ge2M,\qquad a=gkx_0\le a_0,\qquad 0<g<g_0,\quad0<k<k_0,
\end{equation}
without a bound on \(S\). Indeed \(kg/x\le q\) for \(x\ge a\),
so the choice of \(M\) gives
\[
 x+2g\le m_+(x)\le x+4g,\qquad
 6g+a+\pi gq+gq^2\le C(g+a).
\]
These are the two exterior hypotheses with \(\rho=g+a\).
Choose \(d_q\) within that lemma's fixed cut allowance and
\(a_0<d_q/L\), \(L\ge2\), with \(g_0+a_0<\min(\rho_0,1)\).
Shrink \(g_0\) so that \(g_0\delta_q<a_0\).
Its construction gives, for both signs on their respective domains,
\begin{equation}\label{deg:quadratic-outer-domain}
 D\ge c(g+W),\quad 0<W\le w_1,\qquad
 W\ge w_0\quad(x\ge La).
\end{equation}
For the positive sign, in fact \(D\ge c(g+x+W)\).
The limiting graph in \eqref{deg:spectral-exterior-window} keeps
\(t=s\), hence \(m_s^*(x)=sx\). Its positive window is
\[
 |W(d_q)-W_+^*(d_q)|\le C(g+a)+C(a/d_q)^c.
\]
The original regular-section reserves give common neighbourhoods of the
two limiting terminal families whose continuations lie strictly inside
the prescribed raw buffers. Choose \(a_0\), then \(g_0\), inside the
positive reserve. For the negative sign the same window with \(\rho\asymp g\)
and \(a\le g\delta_q\) gives its original buffered image after shrinking
\(g_0\). Thus both regular continuations are evaluated on admitted
images. The positive graph is the inverse of the actual incoming
passage; the negative graph follows positive time, and \(z_x>0\).

For the logarithmic rows on \(S\le\delta_q\), use \(\rho\asymp g\)
for both signs. The coefficient bounds hold on a slightly enlarged open
phase interval, with all coefficients fixed under \(E\) and \(\beta=b_0\).
The admitted regular continuation followed by the observer is the
\(\chi\) of the common lemma; its prescribed ordinary \(C^3\)
logarithmic-slope reserve gives \(\mathcal T\le C\).

Write \(W_l=\ell_0gp\), \(\eta=\log z\), and
\(\nu=6g/(1+2g)=\lambda/2\).
The exact moving factor specializes to
\begin{equation}\label{deg:quadratic-outer-moving-factor}
 \mathscr B_s=EW_l-F_s(a,W_l)
 =\ell_0\left[g(6-s\pi q+3sS)
 -\frac{1+2g+sa+sCa(g+sa)W_l}{1+\ell_0Q_s(a)}\right].
\end{equation}
It is a negative unit. With
\(\mathscr L_s=[2m_s+x m_s'+(xQ_s'+sCx(g+sx))W]/D\),
the exact spectral density is
\begin{equation}\label{deg:quadratic-outer-density}
 \mathscr L_s\,d\tau=\nu\,d\eta+J_s\,d\eta,\qquad
 J_s=\frac{s\{(3-\nu)x-\pi kg^2/x
       +x[-A+(3-\nu)C(g+sx)]W\}}{N_s} .
\end{equation}
Put \(I(a)=\int_{\eta_l(a)}^{\eta_T(a)}J_s\,d\eta\).
For an observer choose a fixed reference \(z_*\) in its buffered image,
put \(c_\theta=\Phi'(z_*)>0\), and \(\widehat\Phi'=\Phi'/c_\theta\).
The logarithmic derivative assumption makes \(\widehat\Phi'\) a positive
unit there. For the raw embedding take \(c_\theta=1\).
The common exact slope, using \(b_0=-a^2/(2g^2k^2)\), is
\begin{equation}\label{deg:quadratic-outer-exact-slope}
 \begin{gathered}
 (\Phi\circ H_s)'(b_0)=d_\theta a^{-\lambda}U_s,
 \quad d_\theta=c_\theta g^2k^2\ell_0^{-1-\nu}g^{-\nu},\\
 U_s=(-\mathscr B_s)(1+\ell_0Q_s(a))p^{-\nu}
       \frac{W_T}{D_T}z_T^\nu\widehat\Phi'(z_T)e^I .
 \end{gathered}
\end{equation}
By \eqref{deg:spectral-exterior-rows}, its first two logarithmic residual
rows satisfy
\[
 |E\log U_s|+|E^2\log U_s|
 \le C\{gL_g(q+S)+a^{2-\lambda}\}\le CgL_g(q+S),
\]
since \(2-\lambda>1\) and \(a=gS\). The spectral factor retains
\(\lambda=12g/(1+2g)\). With \(R_s=E\log U_s\) and \(G_s=h_sR_s\),
\[
 D_s\log\phi_s'=-\lambda h_s+G_s,\qquad
 D_sG_s=h_s(Eh_s)R_s+h_s^2ER_s.
\]
The bounded \(h_s,Eh_s\) give
\eqref{deg:quadratic-outer-weighted-jets}. Separate affine value/slope
calibrations are constant under \(E\) and do not change these rows.
\end{proof}

\begin{corollary}\label{deg:cor-positive-large-S-outer}
Keep the original positive raw section and observer buffers of
Lemma~\ref{deg:lem-quadratic-outer-jets}, and fix \(\delta>0\).
There are \(a_0,g_0,k_0>0\) such that its positive outward passage and
observer are defined on the whole range
\[
 x_0\ge2M,\qquad S=kx_0\ge\delta,\qquad
 a=gS\le a_0,\qquad 0<g<g_0,\quad0<k<k_0.
\]
With the same fixed-coefficient \(E=x_0\partial_{x_0}\), quadratic
label and coefficient-only slope normalization, one has
\begin{equation}\label{deg:positive-large-S-direct-jets}
 |E\log\phi_+'|+|E^2\log\phi_+'|
 \le C_\delta a(1+|\log a|).
\end{equation}
There is no upper bound on \(S\) independent of \(g\), and no comparison
of \(g\) with \(k\). The observer has exactly the same stated
ordinary \(C^3\) logarithmic-slope input as in the lemma.
\end{corollary}

\begin{proof}
Use the common positive admission \eqref{deg:quadratic-positive-admission}
and its original buffered continuation. Apply
Lemma~\ref{deg:lem-spectral-exterior} with the same positive dictionary,
\(\rho=g+a\), \(\sigma=gk\) and terminal cut \(d_q\).
Its coefficient and mass hypotheses were checked before imposing
\(S\ge\delta\); the admitted ordinary \(C^3\) observer again gives
\(\mathcal T\le C\). Now
\[
 g\le a/\delta,\quad
 \Theta=a+\pi gq+gq^2\le C_\delta a,\quad
 L_\rho\le1+|\log a|,\quad
 \lambda=\frac{12g}{1+2g}\le12a/\delta.
\]
With \(e=2-\lambda>3/2\), the spectral rows give
\[
 |E\log\phi_+'+\lambda|+|E^2\log\phi_+'|
 \le C_\delta\{a(1+|\log a|)+(a/d_q)^e\}
 \le C_\delta a(1+|\log a|).
\]
Absorbing \(\lambda=O_\delta(a)\) proves
\eqref{deg:positive-large-S-direct-jets}. The exact quadratic label,
coefficient-only \((gk)^2\) gain and \(a^{-\lambda}\) are retained.
\end{proof}


\begin{corollary}\label{deg:cor-negative-absolute-outer}
Keep the original negative raw section and buffered observers, and fix
\(r_c>0\). For independently small \(g,k>0\), the actual positive-time
negative outward passage is defined for every
\[
 r\ge r_c,\qquad 0<R=kr<\mathfrak c,
 \qquad p=3+\pi/r+r^{-2}-R>0.
\]
Here \(\mathfrak c\) is the negative simple-root endpoint already defined
in the quadratic chart. Put
\[
 E=r\partial_r=R\partial_R,\qquad B=1+g(2-R),\qquad
 h=p/B,\qquad \mathcal A=hE=\mathcal D_-.
\]
With primes in either quadratic label \(-r^2/2\) or \(-R^2/2\),
\begin{equation}\label{deg:negative-absolute-direct-jets}
 |\mathcal A\log\phi_-'|+|\mathcal A^2\log\phi_-'|
 \le Cg(1+|\log g|).
\end{equation}
The ordinary \(C^3\) buffered observer input is the same as in the lemma;
it includes the actual divided finite closing on its original windows.
\end{corollary}

\begin{proof}
Use the exact negative equations of the lemma, with
\(a=gR\), \(q=1/r\), \(W_l=\ell_0gp\).
Here \(p\) may tend to zero at the simple root, so the weak height
interval can be arbitrarily long. The moving initial variation has
a factor \(p\), which will absorb this dwell in the differentiated
integrals. We first construct the graph in logarithmic height.
On the stated range \(p\) is bounded above, \(R\) is bounded,
\(a\le Cg\), and \(m_-(a)=gp\). Also \(m_-'(x)<0\).
To establish the actual graph without presupposing a positive raw
radial denominator, use \(\eta=\log z\), \(\chi=\log x\), and
\[
 \chi_\eta=V=(m_-+Q_-W)/N_-,\qquad W=e^{\eta-2\chi}.
\]
This equation is regular while \(N_-\) is a positive unit.
Initially \(D_l=gp(1+\ell_0Q_-(a))>0\).
At a first zero of \(D=m_-+Q_-W\), its derivative along this graph
would be \(D_\eta=Q_-W>0\). Hence \(D>0\) throughout the graph.
Below a smaller auxiliary positive strip threshold,
\(D\le gp+Q_-W\) and \(V<1/4\); thus \(W_\eta>0\) and
\(W\ge\ell_0gp\). Consequently
\[
 W_\tau=WN_-/D-2W
 \ge\frac{\inf N_-}{\ell_0^{-1}+\sup Q_-}-2W\ge c_0>0.
\]
A fixed logarithmic length therefore reaches that threshold; the fixed
positive margin supplies the remaining climb to the original \(w_0\).
Thus \(x_1\le La\), uniformly even when \(p\) tends to zero.
Choose \(g_0\) so that \(La<d_q\). Thereafter the original strip is
invariant and \(D\ge cW\). The limiting strong graph and contraction
estimate of the lemma give
\[
 |W(d_q)-W_-^*(d_q)|\le Cg+C(La/d_q)^c\longrightarrow0.
\]
This proves first passage and admission to the original observer buffer
before derivative estimates. Its initial height length satisfies
\[
 \eta_1-\eta_l
 =\log\frac{W_1}{\ell_0gp}+2\log\frac{x_1}{a}
 \le C(1+|\log g|+|\log p|).
\]

Write \(Q_-=q_0+q_1x+Cx^2\), where
\(q_0=1-Ag+Cg^2\), \(q_1=A-2Cg\).
At fixed \(\eta\) the exact partial derivatives are
\[
 D_\chi=-x-\pi kg^2/x-2k^2g^3/x^2-(2q_0+q_1x)W,
 \qquad N_{-,\chi}=-x+CgxW.
\]
On a sufficiently small fixed inner raw interval they imply
\begin{equation}\label{deg:negative-inverse-dissipation}
 V_\chi\le-c\Xi,\qquad |V_{\chi\chi}|\le C\Xi,
 \qquad \Xi=x+kg^2/x+k^2g^3/x^2+W.
\end{equation}
Indeed \(D\le gp+Q_-W\le Cg+CW\); the correction
\(-DN_{-,\chi}/N_-^2\) is absorbed by the negative terms in
\(D_\chi/N_-\). The next derivative has the same density.

Let \(X_j=\mathcal A^j\chi|_\eta\),
\(l_j=\mathcal A^j\eta_l\), and \(t_j=\mathcal A^j\eta_T\).
The full moving initial rows are
\begin{align*}
 l_1&=(2p+Ep)/B,\qquad l_2=O(p),\\
 X_{1,l}&=h-V_ll_1=pK_l>0,\\
 X_{2,l}&=\mathcal A X_{1,l}-V_{\chi,l}X_{1,l}l_1=O(p).
\end{align*}
with \(K_l\) a positive unit. Here \(l_1\), \(Eh\) and \(El_1\)
are bounded. The two inverse equations are \eqref{deg:scalar-inverse-endpoints}.
Up to that inner cut, \eqref{deg:negative-inverse-dissipation} gives
\(0<X_1\le Cp\), \(\int\Xi X_1\,d\eta\le Cp\), and
\(|X_2|\le CX_1\), by division by the positive first variation.
On the strong part choose an auxiliary fixed inner cut \(d_*\le d_q\)
so that \(V_\chi/V\le-3/2\), then shrink \(g_0\) to ensure
\(La<d_*\). The decay follows from
\(D\ge cW\), \(W\ge w_0\), and
\(V_\chi/V\le-2+Cg/w_0+Cx\).
Thus \(X_1\le Cp(a/x)^{3/2}\) up to that cut.
The remaining compact strong-graph segment to the original \(d_q\)
has bounded coefficients and finite logarithmic length, and preserves
this bound up to a fixed factor. At the original terminal cut, \eqref{deg:scalar-inverse-endpoints}
gives \(|t_1|+|t_2|\le Cp a^{3/2}\).
Use the exact density \(J_-=J\) from
\eqref{deg:quadratic-outer-density}, with
\(I=\int_{\eta_l}^{\eta_T}J\,d\eta\).
Its two complete rows are \eqref{deg:scalar-density-endpoints},
with \(D=\mathcal A\), \(r=\chi\), and \(Dr_l=h\).
On the initial part \(|\partial_\chi^jJ|\le Cg\), \(j=0,1,2\),
since \(x\le La\le Cg\) and \(kg^2/x\le g/r\).
Its differentiated integrals are at most
\(Cgp(1+|\log g|+|\log p|)\le Cg(1+|\log g|)\).
On the strong part \(d\eta\le C\,d\log x\), and the integrals are
bounded by
\[
 Cp\int_{La}^{d_*}(x+kg^2/x)(a/x)^{3/2}\,d\log x
 \le Cp(a+kg^2/a)\le Cg.
\]
If the actual entry is below \(La\), the fixed logarithmic segment
from \(x_1\) to \(La\) has the same \(Cg\) bound.
The compact final segment contributes \(Cp a^{3/2}\).
The displayed initial terms are \(O(g)\), and all terminal partials
are bounded. Therefore \(|\mathcal AI|+|\mathcal A^2I|\le CgL_g\).

Finally use the exact complete slope factor of the lemma.
Its moving-initial and \(1+\ell_0Q_-(a)\) logarithmic rows are
\(O(g)\). Although \(p^{-\nu}\) is unbounded as \(p\) tends to zero,
its two action logarithmic rows are \(O(g)\), since
\(\mathcal A\log p=Ep/B\) and \(\mathcal A^2\log p=O(p)\).
The terminal \(W_T/D_T\), \(z_T^\nu\) and ordinary \(C^3\) observer
rows are \(O(p a^{3/2})\), by the full terminal formulas.
The extracted \(a^{-\lambda}\) contributes \(-\lambda h=O(g)\)
and \(-\lambda\mathcal Ah=O(g)\).
All coefficient-only slope factors are frozen, including the change
between the two quadratic labels. These estimates prove
\eqref{deg:negative-absolute-direct-jets} on the whole original domain,
including arbitrarily close approach to its negative root endpoint.
\end{proof}

\subsection{The complete negative inward action remainder}

We next continue the negative lip back to the original centre and
estimate its action remainder. This supplies one side of the later
matching equation. The positive and negative heights are constructed
separately; matching is imposed only on the intersection of their images.

\begin{corollary}\label{deg:cor-negative-inward-common-centre}
Use \eqref{deg:quadratic-inward-family} with \(s=-1,x_s=r\),
\(\pi\in[\pi_\flat,\Pi]\), and fixed small \(\ell_0\).
Fix \(r_c>0\). On the entire strict negative interval
\(r\ge r_c\), \(0<R=kr<\mathfrak c\), continue the lip
\(Y=\ell_0P_-(r)\) to \(\xi=0\) along the inverse of the actual
positive-time negative branch. Its height \(Z_-(r)>0\) is defined on
this entire interval and enters the original central height window
after one coefficient shrink. In \eqref{deg:quadratic-inward-memory-common}, its remainder obeys
\begin{equation}\label{deg:negative-inward-central-remainder}
 \sum_{j=1}^3|\mathcal D_-^j\mathcal R_-|\le Cg.
\end{equation}
Changing the central height by a coefficient-only multiplicative scale changes
only a constant in this logarithmic formula. The estimate does not
differentiate coefficients or bound the value of \(\mathcal R_-\) by
\(Cg\). All constants are independent of the distance to the negative
simple root and of any comparison between \(g\) and \(k\).
\end{corollary}
\begin{proof}
Choose an auxiliary fixed \(0<b<r_c/2\). On \([0,r_c]\), the
polynomial is uniformly positive after decreasing \(k\), and
\[
 A_-(r_c)-A_-(b)\ge a_*>0 .
\]
For \(v\ge b\), \(kv<\mathfrak c\), write
\[
 Q=1-Ag(1-kv)+Cg^2(1-kv)^2,\quad
 H=\frac{P_-'(v)}{vB_-(v)},\quad
 \omega=-\frac{Ckg^3(1-kv)P_-(v)}{vB_-(v)} .
\]
Here \(Q\) is a positive unit. With
\(\mathcal D_v=P_-(v)\partial_v/(vB_-(v))\), direct rational
differentiation in \((v^{-1},kv)\) gives
\begin{equation}\label{deg:negative-inward-coefficient-rows}
 \sum_{j=0}^3|\mathcal D_v^jH|\le C,\qquad
 \sum_{j=1}^3|\mathcal D_v^jQ|\le Cg,\qquad
 \sum_{j=0}^3|\mathcal D_v^j\omega|\le Cg^3 .
\end{equation}
The field \(\mathcal D_v\) contains the factor \(P_-(v)\);
none of these bounds divides by its vanishing root factor.

Set \(t=A_-(r)\), \(q=(t-A_-(v))/g\), and
\(W=Y/P_-(v)\). The original negative scalar equation gives exactly
\[
 v_q=-g\frac{P_-(v)}{vB_-(v)},\qquad
 (\log W)_q=-\frac{1+\omega W}{1+QW}+gH,
 \qquad W(0)=\ell_0 .
\]
At fixed \(q\), \(\partial_t v=P_-(v)/(vB_-(v))\), so every
\(\partial_t\) coefficient derivative is precisely \(\mathcal D_v\).
The logarithmic rate and the normal diagonal are at most \(-c<0\)
for the fixed small lip. Lemma~\ref{deg:lem-stable-ratio-ledger}
applies with \(D=\partial_t\), \(\kappa=g\), \(\rho=g\),
\(\alpha=0\), by \eqref{deg:negative-inward-coefficient-rows};
clock differentiation acts as \(-g\mathcal D_v\). It constructs the
actual positive-height inverse passage to \(b\), with radial speed
\(P_-+QY>0\), and gives
\[
 |\partial_t^jW|+|\partial_t^jL|
 \le Cg(1+q)^Ne^{-cq},\quad1\le j\le3,\qquad
 L=(Q-\omega)W/(1+QW).
\]
Retain the exact source polynomial:
\[
 T=(t-A_-(b))/g,\qquad V_b=\int_0^{T(t)}L\,dq,\qquad
 Y_b=\ell_0P_-(r)e^{-T+V_b}=P_-(b)W(T)\le Ce^{-cT}.
\]
Here \(T'=1/g,T''=T'''=0\), and \(T\ge a_*/g\).
The three complete rows \eqref{deg:stable-ratio-terminal-rows}
therefore give \(|\partial_t^jV_b|\le Cg\): each boundary term
is bounded by \(Cg^{-N}(1+T)^Ne^{-cT}\). The retained polynomial
gives
\[
 \mathcal D_-\log Y_b=H(r)-g^{-1}+\partial_tV_b,
 \qquad |\mathcal D_-^jY_b|\le Cg^{-j}Y_b,\quad j\le3.
\]

Reflect the segment from \(-b\) to the original centre by \(v=-\xi\).
Preparation~\eqref{deg:positive-large-S-inward-compact-graph}, with \(s=-1\),
then applies on the fixed margin above. Exponential flatness gives
\(Y_b\le h_b/2\) and small \(Y_b/g\), so
Lemma~\ref{deg:lem-compact-height-continuation} establishes the actual
continuation, \(0<Y\le2Y_b\), and its logarithmic nonlinear unit
\(\mathcal N\). The three composition rows of that lemma are
bounded by \(CY_b/g^{j+1}\le Cg\).
The exact linear multiplier is \(e^{-A_-(b)/g}\), so the artificial
cut cancels and
\(\mathcal R_-=V_b+\mathcal N(Y_b)\) gives
\eqref{deg:negative-inward-central-remainder} at the original centre.
\end{proof}

\subsection{Positive inward passage on the growing range}

When \(S=ku\) grows, the positive action speed is of order
\(1+S\). We therefore estimate the inward remainder in the radial
derivative \(u\partial_u\), keeping the conversion to action
derivatives explicit. This is the estimate used in
\eqref{deg:quadratic-large-S-curvature}.

For \(s=+1\), the factorization is
\[
 P(v)=(\mathfrak c+kv)\{(v-d)^2-\Delta\},\quad
 \mathfrak c=3+2kd,\quad \mathfrak c(d^2-\Delta)=1,\quad
 \pi=2\mathfrak cd-k/\mathfrak c,
\]
with \(d=\pi/6+O(k)\), \(\mathfrak c=3+O(k)\).

\begin{corollary}[Positive large-\(S\) inward action rows]
\label{deg:cor-positive-large-S-inward}
Fix \(0<\pi_\flat\le\Pi<\infty\), a compact set of the passive
coefficients \(A,C\), and \(\delta>0\). Choose the original lip
height \(\ell_0\) and physical collar width \(d_0\) sufficiently small
within their original buffered windows. There are fixed
\[
 v_i,M,g_0,k_0>0,\qquad
 k_0\le\delta/(2M),
\]
such that the following holds for independently small
\(0<g<g_0\), \(0<k<k_0\), \(\pi\in[\pi_\flat,\Pi]\).
On every strict root-free fibre \(\Delta<0\), every lip point
\begin{equation}\label{deg:positive-large-S-inward-domain}
 \xi=u,\qquad Y=\ell_0P(u),\qquad
 S=ku\ge\delta,\qquad a=gS<d_0
\end{equation}
has an actual positive-height, positive-time first passage to \(\xi=0\).
Its original centre height is \(Y_{0,+}=Z_+>0\), with
\(\mathcal R_+\) defined by \eqref{deg:quadratic-inward-memory-common}.
With all coefficients fixed under \({\mathcal E}=u\partial_u\),
\begin{equation}\label{deg:positive-large-S-inward-action3}
 |{\mathcal E}^j\mathcal R_+(u)|\le C_\delta a,\qquad j=1,2,3 .
\end{equation}
The constants are independent of \(g,k\) and of \(-\Delta>0\).
The statement can be restricted to the original strict raw-window
preimages without extending any image or inverse.

On a real or double-root fibre the prescribed positive through lip ends
at the first root \(a_{\rm root}=d-\sqrt{\Delta}\). That endpoint is
bounded uniformly in this coefficient block, and after the same
\(k\)-shrink its through interval has \(ku<\delta\). Thus
\eqref{deg:positive-large-S-inward-domain} is empty on that branch. No passage beyond the first
root, or at the stopped \(\Delta=0\) limit, is included.
\end{corollary}

\begin{proof}
Fix the finite differentiability budget used below first.
Choose \(v_i>0\) so small that
\[
 P(v)\ge p_i>0\qquad(0\le v\le2v_i)
\]
uniformly for the fixed \(\pi\)-interval and small \(k\).
Then fix \(M>2v_i\) beyond all positive roots and turning points of
the retained compact positive factors, and large enough that
\begin{equation}\label{deg:positive-large-S-inward-growing-polynomial}
 P(v)\asymp v^2(3+kv),\qquad v\ge M .
\end{equation}
The constants may depend on the fixed passive compact,
\(\pi_\flat,\Pi,\delta,\ell_0,v_i,M\).
All calculations are also valid on a slightly enlarged phase range
\(S>\delta/2\), \(a<2d_0\); the choices of \(d_0,g_0,k_0\) are made
for that range before restriction to \eqref{deg:positive-large-S-inward-domain}.
This supplies actual phase buffers for the displayed derivatives.

\smallskip
\noindent\emph{Actual passage and value admission.}
For \(\Delta<0\), \(P>0\) on the positive half-axis.
On \(v_i\le v\le u\) put
\begin{equation}\label{deg:positive-large-S-inward-action-variable}
 q=\frac{J(u)-J(v)}g,\qquad
 v_q=-g\frac{P(v)}{vB(v)},\qquad
 W=\frac{Y}{P(v)},\qquad
 H=\frac{P'(v)}{vB(v)},\quad
 \varpi=\frac{\zeta(v)P(v)}{vB(v)} .
\end{equation}
Division of \eqref{deg:quadratic-inward-scalar-common}, retaining the polynomial factor, gives
exactly
\begin{equation}\label{deg:positive-large-S-inward-ratio}
 W_q=W\left\{-\frac{1+\varpi W}{1+QW}+gH\right\},
 \qquad W(0)=\ell_0 .
\end{equation}
Here \(B,Q\) are positive units, \(gH\le C_\delta a\), and
\(|\varpi|\le C_\delta a^3\), as also checked below.
Choose \(\ell_0,d_0,g_0\) so that the logarithmic rate in
\eqref{deg:positive-large-S-inward-ratio} is at most \(-c_0<0\) while \(0<W\le\ell_0\).
A first-exit argument proves
\begin{equation}\label{deg:positive-large-S-inward-value-decay}
 0<W(q)\le\ell_0e^{-c_0q},\qquad
 0\le q\le T(u):=\frac{J(u)-J(v_i)}g .
\end{equation}
Consequently \(P+QY=P(1+QW)>0\). For each strict ghost coefficient
the finite interval \([v_i,u]\) has a positive minimum of \(P\);
the time integral to \(v_i\) is finite. Height is positive by
uniqueness. The same construction continues to \(v_i/2\), on its
fixed positive-\(P\) buffer, before imposing the cut \(v_i\).
Thus the fixed-\(q\) partials at that moving terminal are genuine
nearby actual-flow derivatives.

Since \(Y\le\ell_0P(v)\), the physical height on the entire arc is
at most
\[
 C\ell_0\{g^3k^2+g^3k^2v+g^3k^2v^2+g^3k^3v^3\}
 \le C\ell_0\{g^3k^2+g^2ka+ga^2+a^3\}.
\]
Together with \(0<gkv\le a<d_0\), this admits the arc to a
fixed small original chart. Neither a limiting axis orbit nor a
normalization is used to infer this admission.

On \(M\le v\le\delta/k\), \(P(v)\le C_\delta v^2\) and \(B\ge1\).
Since \(u\ge\delta/k\), it follows that
\begin{equation}\label{deg:positive-large-S-inward-action-gap}
 T(u)\ge\frac{c_\delta\log(1/k)-C_\delta}{g}.
\end{equation}
Choose \(k_0\) so that the numerator has a fixed positive margin.
Then \(T\) is uniformly large for small \(g\), with no comparison
between \(g\) and \(k\). In particular the actual terminal height
\begin{equation}\label{deg:positive-large-S-inward-inner-value}
 y_i(u):=Y(v_i;u)=P(v_i)W(T(u),u)
              \le C e^{-c_0T(u)}
\end{equation}
lies in any fixed sufficiently small compact-entry interval
\(0<y_i<y_*\) after the permitted coefficient shrink.
This places the input in the prescribed compact-continuation buffer
before its derivatives are estimated.

On \([0,v_i]\), preparation~\eqref{deg:positive-large-S-inward-compact-graph}
and exponential flatness give the prescribed height buffer and small
\(y_i/g\). Lemma~\ref{deg:lem-compact-height-continuation} therefore
continues the actual graph to the original centre with \(0<Y\le2y_i\),
for either sign of \(C\).
Retain the exact memory~\eqref{deg:stable-ratio-polynomial-memory},
with \(V(q,u)=\int_0^qL\,ds\), \(L=(Q-\varpi)W/(1+QW)\);
its initial factor is \(P(u)/P(v)\) and \(|V|\le C\ell_0\).

\smallskip
\noindent\emph{Three transported coefficient rows.}
Write
\[
 h(v)=\frac{P(v)}{v^2B(v)},\qquad
 {\mathcal D}_v=\frac{P(v)}{vB(v)}\partial_v=h(v)v\partial_v,\qquad
 b(u)=\frac1{h(u)} .
\]
At fixed elapsed action \(q\), differentiating
\(J(v)=J(u)-gq\) yields
\begin{equation}\label{deg:positive-large-S-inward-transport}
 {\mathcal E}\log v=\frac{h(v)}{h(u)},\qquad
 {\mathcal E}[F(v(q,u))]=b{\mathcal D}_vF .
\end{equation}
The full second and third transported rows are
\begin{equation}\label{deg:positive-large-S-inward-transport3}
\begin{split}
 {\mathcal E}^2F&=b^2{\mathcal D}_v^2F+({\mathcal E} b){\mathcal D}_vF,\\
 {\mathcal E}^3F&=b^3{\mathcal D}_v^3F+3b({\mathcal E} b){\mathcal D}_v^2F+({\mathcal E}^2b){\mathcal D}_vF .
\end{split}
\end{equation}
The variable coefficients of \({\mathcal D}_v\) are differentiated in its
powers. On the initial domain,
\begin{equation}\label{deg:positive-large-S-inward-initial-speed}
 h(u)=\frac{3-\pi/u+u^{-2}+S}{1+2g+gS}
       \asymp_\delta1+S,\qquad
 |{\mathcal E}^jh(u)|\le C_\delta h(u)\quad(j=1,2).
\end{equation}
Hence
\[
 |{\mathcal E} b|+|{\mathcal E}^2b|\le C_\delta b,\qquad
 b\le C_\delta/(1+S),\qquad
 gh(u)\asymp_\delta g+a .
\]

For \(v\ge M\), let \(t=1/v\), \(s=kv\), and
\({\mathcal E}_v=-t\partial_t+s\partial_s\). Then
\begin{equation}\label{deg:positive-large-S-inward-rational-symbols}
 h(v)=\frac{3-\pi t+t^2+s}{1+2g+gs},\qquad
 H(v)=\frac{6-\pi t+3s}{1+2g+gs},\qquad
 \varpi(v)=Cg^3s(1+s)h(v).
\end{equation}
For \(0\le s\le S\), \(t\le1/M\), and \(gS=a<2d_0\),
these rational symbols give
\[
 h(v)\le C_\delta h(u),\qquad
 |{\mathcal E}_v^jh(v)|\le C_\delta h(v)\quad(j\le2).
\]
The first three \({\mathcal E}_v\) rows of \(gH,Q-1,\varpi\) have
size \(C_\delta a\); the \(\varpi\) rows have the sharper
size \(C_\delta a^3\). Expanding \((h(v){\mathcal E}_v)^j\) and
using \(b h(v)\le C_\delta\) proves the same bounds for
the rows in \eqref{deg:positive-large-S-inward-transport3}.

On \(v_i/2\le v\le M\), \(H\) has uniform ordinary finite
derivatives. The coefficients and finite ordinary derivatives of
\({\mathcal D}_v=P/(vB)\partial_v\) are uniformly bounded, also as
\(-\Delta\downarrow0\); they contain \(P\) rather than \(P^{-1}\)
at the bottleneck. The corresponding ordinary rows of \(gH,Q-1\)
are \(O(g)\), and those of \(\varpi\) are \(O(kg^3)\).
Since \(g\le a/\delta\) on the retained domain, this proves
\begin{equation}\label{deg:positive-large-S-inward-coefficient-Euler}
 |{\mathcal E}^j(gH)|+|{\mathcal E}^j(Q-1)|+|{\mathcal E}^j\varpi|
      \le C_\delta a,\qquad 0\le j\le3,
\end{equation}
uniformly through the arbitrary strict ghost gap.
In particular the quotient in \eqref{deg:positive-large-S-inward-transport} and its first
two Euler rows are bounded on both the growing and compact portions.

For the time rows, \(\partial_q=-g{\mathcal D}_v\) on coefficient
functions and \(gh(v)\le C_\delta a\) on the growing portion.
On the compact portion \(g{\mathcal D}_v\) has bounded finite coefficients.
The same rational calculations give bounded first and second
\(q\)-rows and the mixed \(q,{\mathcal E}\) rows required below.
Their rows containing an \({\mathcal E}\) retain the \(C_\delta a\) factor.
No inverse of the bottleneck polynomial is introduced.

\smallskip
\noindent\emph{Normal variations and the moving terminal.}
The exact diagonal is
\begin{equation}\label{deg:positive-large-S-inward-negative-diagonal}
 \mathscr F_W=-\frac{1+2\varpi W+\varpi QW^2}{(1+QW)^2}+gH
 \le-c_1<0.
\end{equation}
Thus Lemma~\ref{deg:lem-stable-ratio-ledger} applies with
\(D=\mathcal E,\kappa=g,\rho=C_\delta a,\alpha=0\), by
\eqref{deg:positive-large-S-inward-coefficient-Euler} and the mixed
clock bounds above. Its fixed-clock estimate is
\begin{equation}\label{deg:positive-large-S-inward-W-Euler}
 |\mathcal E^jW|+|\mathcal E^jL|
 \le C_\delta a(1+q)^Ne^{-cq},\qquad1\le j\le3,
\end{equation}
with the mixed and pure clock rows specified there. The terminal rows are
\begin{equation}\label{deg:positive-large-S-inward-terminal-clock}
 T_1=b/g,\qquad T_2=(\mathcal E b)/g,\qquad
 T_3=(\mathcal E^2b)/g,\qquad
 |T_j|\le C_\delta/(gh(u))\le C_\delta/a\le C_\delta/g .
\end{equation}
By \eqref{deg:positive-large-S-inward-action-gap}, for each fixed \(N\),
\begin{equation}\label{deg:positive-large-S-inward-terminal-absorption}
 g^{-N}(1+T)^Ne^{-cT}\le C_{\delta,N}g\le C_{\delta,N}a.
\end{equation}
Choose \(k_0\) with \(cc_\delta\log(1/k_0)-cC_\delta>0\);
then \(g^{-N}e^{-\beta/g}\) absorbs every terminal loss, using
independent upper bounds on \(g,k\).
The terminal formulas~\eqref{deg:stable-ratio-terminal-rows}, for
\(V_i=V(T(u),u)\), give
\begin{equation}\label{deg:positive-large-S-inward-V-action3}
 |\mathcal E^jV_i|\le C_\delta a,\qquad1\le j\le3.
\end{equation}
The exact memory, including all moving terminal rows, is
\begin{equation}\label{deg:positive-large-S-inward-height-memory}
 y_i=\ell_0P(u)e^{-T+V_i}=P(v_i)W(T),\qquad
 \mathcal E^j\log y_i=\mathcal E^j\log P(u)-T_j+\mathcal E^jV_i.
\end{equation}
The rational symbols bound \(|\mathcal E^j\log P(u)|\le C_\delta\),
hence \(|\mathcal E^j\log y_i|\le C_\delta/g\) for \(1\le j\le3\).
The value estimate uses \(P(v_i)W(T)\), independently of \(P(u)\).
The compact continuation already constructed has exact multiplier
\(m(0)=e^{-J(v_i)/g}\). Its three composition rows from
Lemma~\ref{deg:lem-compact-height-continuation} are at most
\(Cy_i/g^4\le C_\delta a\), by
\eqref{deg:positive-large-S-inward-terminal-absorption}. Consequently
\[
 \log Y_{0,+}=\log\{\ell_0P(u)\}-J(u)/g+V_i+\mathcal N(y_i),
 \qquad \mathcal R_+=V_i+\mathcal N(y_i).
\]
This proves the three Euler bounds at the original centre. The fixed-cut
action cancels exactly; an additive action normalization changes only a
coefficient constant in \(\mathcal R_+\).

Finally, \(d,\mathfrak c\), and every nonnegative
\(\sqrt{\Delta}\) in the fixed \(\pi\)-block are bounded for
small \(k\). The root \(a_{\rm root}\) is therefore bounded by
the fixed \(M\). Its prescribed through lip has
\(ku\le ka_{\rm root}<\delta\) after the \(k\)-shrink, as stated.
The strict ghost proof permits arbitrarily small \(-\Delta>0\)
but does not add a through interval at \(\Delta=0\).
\end{proof}

\subsection{Bounded positive quadratic passages}

On a bounded positive lip interval, both the root gap and the lip's
distance to its first root can vanish. The inward calculation remains
uniform in action coordinates. For the outward calculation we use
logarithmic height through the weak stage; its moving initial factor
absorbs the long dwell before the graph enters the parabolic strip.

\begin{corollary}[Bounded positive quadratic suppliers]
\label{deg:cor-positive-compact-suppliers}
Fix the original positive raw section and its buffered regular observer,
with its original ordinary \(C^3\) logarithmic-slope reserve
\(\|\partial_z\log\Phi'\|_{C^1}\le C\) on the buffered raw image,
a compact neighborhood of the passive coefficients \(A,C\), and
\(\pi\in[\pi_\flat,\Pi]\). Fix \(U<\infty\) and use
\eqref{deg:quadratic-inward-family}--\eqref{deg:quadratic-inward-scalar-common},
with physical coordinates \(x=gkv,z=g^3k^2Y\).
Retain the original actual moving-lip interval \(Y=\ell_0P(u)\),
intersected with \(u\le U\). Its lower threshold satisfies
\(u_{\rm th}\ge u_0>0\). On real or fold fibres this interval has
\(u_{\rm th}\le u<a\), where \(a\) is the first positive root.
On root-free fibres it is the component of the original section-image
preimage specified in the positive quadratic domain construction.
In each case \(P>0\) on the actual inward axis path \([0,u]\).
No right-hand axis continuation is added at the fold.

There are \(g_0,k_0,C_0>0\), uniform on these domains, such that
the following statements hold for independently small
\(0<g<g_0\), \(0<k<k_0\). Put \(L_g=1+|\log g|\), with \(J,\mathcal D_+\) as in
\eqref{deg:quadratic-inward-memory-common}.
The inward lip orbit reaches the original centre \(v=0\), and its
height \(Z_+(u)=Y(0;u)>0\) has
\begin{equation}\label{deg:positive-compact-central-jets}
 |\mathcal D_+^j\mathcal R_+|\le C_0g,\qquad1\le j\le3.
\end{equation}
The outward lip graph, which is the inverse of the physical positive
incoming passage, reaches the original positive raw port at positive
height and inside its original buffer. For the complete embedding
followed by the original increasing ordinary \(C^3\) observer, use
the lip label \(\beta=-u^2/2\) and any coefficient-only positive
affine value and slope calibration \(\phi_+\). Then
\begin{equation}
 \phi_+'>0,\qquad
 |\mathcal D_+^j\log\phi_+'|\le C_0gL_g,\quad j=1,2.
 \label{deg:positive-compact-outer-jets}
\end{equation}
The prime denotes differentiation in \(\beta\).
Both estimates freeze \(g,k,\pi,A,C\). They are uniform as the
positive root gap and the distance of the lip to its first root tend
to zero independently of \(g,k\). They do not require an intersection
of positive and negative inward images.
\end{corollary}

\begin{proof}
The exact positive quadratic factorization is
\[
 P(v)=(\mathfrak c+kv)\{(v-d)^2-\Delta\},\qquad
 \mathfrak c=3+2kd,\quad
 \mathfrak c(d^2-\Delta)=1,\quad
 \pi=2\mathfrak cd-k/\mathfrak c.
\]
It gives uniform bounds for
\(d,d^{-1},\mathfrak c,\mathfrak c^{-1}\) on the retained fold band.
On a real fibre, \(P(a)=0\) implies
\[
 1=\pi a-3a^2-ka^3\le\Pi a,\qquad a\ge\Pi^{-1}.
\]
On a root-free fibre,
\(\pi=2\mathfrak cd-k/\mathfrak c\) bounds \(d\) away from zero.
Thus the original fixed \(X_{\rm up}\) can be chosen small enough
that \(u_{\rm th}=a-X_{\rm up}\), respectively \(d-X_{\rm up}\),
has the asserted uniform positive lower bound. All constants below
are fixed before \(g,k\).

\smallskip
\noindent\emph{The inward action integral and compact continuation.}
Choose fixed \(b<u_0/2\) with \(P\ge p_b>0\) on \([0,b]\), and put
\[
 \begin{gathered}
 q=(J(u)-J(v))/g,\quad W=Y/P(v),\quad H=P'/(vB),\\
 \omega=\zeta P/(vB),\qquad
 \mathcal D_v=P(v)\partial_v/(vB(v)).
 \end{gathered}
\]
Use the ratio equation and exact memory
\eqref{deg:stable-ratio-polynomial-memory}, with \(\kappa=g\)
and \(V(q)=\int_0^q(Q-\omega)W/(1+QW)\,ds\).
At fixed \(q\), \(\mathcal D_+v=P(v)/(vB(v))\), so its repeated
coefficient rows are powers of \(\mathcal D_v\). On \([b,U]\),
this field has bounded finite derivatives, uniformly at the root/fold:
\begin{equation}\label{deg:positive-compact-inward-coefficients}
 |\mathcal D_v^j(gH)|+|\mathcal D_v^j(Q-1)|+|\mathcal D_v^j\omega|
 \le Cg,\qquad0\le j\le3.
\end{equation}
Clock rows use \(-g\mathcal D_v\), and the fixed small lip gives
both negative margins. These are the inputs of the negative inward proof,
now with \(s=+1\), \(D=\mathcal D_+\), \(\kappa=g,\rho=g,\alpha=0\).
The terminal check is
\[
 q_b=(J(u)-J(b))/g\ge\eta/g,\qquad
 \mathcal D_+q_b=1/g,\quad \mathcal D_+^2q_b=\mathcal D_+^3q_b=0;
\]
the lower bound follows on a fixed subinterval of \([b,u_0]\).
Lemma~\ref{deg:lem-stable-ratio-ledger} and its terminal formulas give
\(|\mathcal D_+^jV(q_b)|\le Cg\), \(1\le j\le3\), and
\begin{equation}\label{deg:positive-compact-central-incoming-height}
 \begin{gathered}
 y_b=\ell_0P(u)e^{-q_b+V(q_b)}=P(b)W(q_b)\le Ce^{-c\eta/g},\\
 \mathcal D_+\log y_b=-g^{-1}+H(u)+\mathcal D_+V(q_b),\qquad
 |\mathcal D_+^jy_b|\le C_jg^{-j}y_b\quad(1\le j\le3).
 \end{gathered}
\end{equation}
Preparation~\eqref{deg:positive-large-S-inward-compact-graph} verifies
the compact-continuation inputs: \(P\ge p_b\),
\(|\partial_Y^jK|\le C/g\), \(y_b\le h_b/2\), and small \(y_b/g\).
Lemma~\ref{deg:lem-compact-height-continuation} then gives
\[
 Z_+(u)=e^{-J(b)/g}y_be^{\mathcal N(y_b)},\qquad
 \mathcal R_+=V(q_b)+\mathcal N(y_b).
\]
Its three composition rows are \(O(y_b/g^{j+1})=O(g)\).
The cut action cancels, proving~\eqref{deg:positive-compact-central-jets}.

\smallskip
\noindent\emph{Construction of the outward graph and buffer admission.}
Put
\[
 a=gku,\qquad p=\frac{P(u)}{u^2}>0,\qquad
 E=u\partial_u,\qquad h=\frac p{B(u)},\qquad \mathcal D_+=hE.
\]
In the raw graph set \(W=z/x^2\), \(r=\log x\), and \(\eta=\log z\).
Its exact coefficients are
\[
 \begin{split}
 m(x)&=3g+x-\pi kg^2/x+k^2g^3/x^2,\\
 Q_0(x)&=1-A(g+x)+C(g+x)^2,\\
 N(x,W)&=1+2g+x+Cx(g+x)W,\qquad D_0=m+Q_0W.
 \end{split}
\]
The initial data are the original lip data
\begin{equation}
 x_l=a,\qquad W_l=\ell_0gp,\qquad
 z_l=a^2\ell_0gp=g^3k^2\ell_0P(u).
 \label{deg:positive-compact-raw-initial}
\end{equation}
Use the inverse logarithmic-height equation first:
\begin{equation}
 r_\eta=V=\frac{D_0}N,\qquad W_\eta=W(1-2V).
 \label{deg:positive-compact-inverse-height}
\end{equation}
Here \(N,Q_0\) are positive units on the original small corridor.
At a possible first zero of \(D_0\), one has \(V=0\) and
\[
 \frac{dD_0}{d\eta}=Q_0W>0.
\]
Thus a graph starting with \(D_0>0\) cannot cross to \(D_0\le0\).
Both \(x\) and \(z\) increase on this graph, including at an axis
root crossed at positive height.

Fix \(L=2\), then a sufficiently small auxiliary height \(w_0>0\).
In the box \(a\le x\le La\), \(0<W\le w_0\), the variable
\(x/(gk)\) lies in the fixed interval \([u_0,LU]\). Hence
\[
 |m|+|m_r|+|m_{rr}|\le Cg,\qquad
 0<V\le C(g+w_0),\qquad \tfrac12W\le W_\eta\le W.
\]
Write \(s=\eta-\eta_l\). The length of this weak stage, up to either
possible first exit, is bounded by
\begin{equation}
 T\le C\{L_g+\log^+(1/p)\}.
 \label{deg:positive-compact-weak-length}
\end{equation}
Since \(m(0)=gp\) and
\[
 |m_s|\le Cg(|m|+W),
\]
Gronwall's inequality, together with the lower growth rate for \(W\),
gives
\[
 |m(s)|\le Cgp e^{Cgs}+CgW(s),\qquad
 r(s)-r_l\le Cgp(1+s)e^{Cgs}+Cw_0.
\]
For every fixed \(j\), the elementary bound
\begin{equation}
 p(1+T)^j e^{CgT}\le C_jL_g^j,\qquad 0<p\le p_{\max},
 \label{deg:positive-compact-polynomial-absorption}
\end{equation}
holds after shrinking \(g\).
Indeed \eqref{deg:positive-compact-weak-length} reduces it to
\(g^{-Cg}\le C\) and the boundedness of
\(p^{1-Cg}\{1+\log^+(1/p)\}^j\).
Choose \(w_0\) so \(Cw_0<\frac12\log L\), then \(g\) so that
\(CgL_g<\frac12\log L\).
The preceding displacement estimate excludes an exit at \(x=La\).
Thus the graph reaches \(W=w_0\) at a finite height time and an
abscissa \(a_*\in[a,La]\), for every \(p>0\).

Keep the original raw cut \(d_q\) fixed. Choose an auxiliary cut
\(0<d_*<d_q\) sufficiently small compared with \(w_0\), before
shrinking \(g,k\). After this shrink \(a_*<d_*/2\).
On \(a_*\le x\le d_*\) use
\[
 W_{\log x}=F=W(N/D_0-2).
\]
The strip \(w_0\le W\le w_1\), for a fixed \(w_1>1/2\), is
invariant, with \(D_0\ge cW\).
Moreover \(F_W\le-c\) and \(V_r/V\le-\alpha\), where
\(\alpha=3/2\), because
\(|m|+|m_r|\le C(g+x)\) on this segment and \(g,d_*\) are
sufficiently small compared with \(w_0\).
The limiting strong graph at \(g=k=0\) has
\(W^*(x)=1/2+O(x)\): with \(U_*=W^*-1/2\) its equation is
\(xU_*'=-2U_*+O(x)\), whose bounded solution follows from the
integral equation at zero.
Since \(x\ge a\ge gku_0\), the strip equation differs by \(O(g)\);
variation of constants gives
\[
 |W(d_*)-W^*(d_*)|\le Cg+C(a_*/d_*)^c.
\]
On the remaining fixed segment \([d_*,d_q]\),
\(F=1-2W+O(g+x)\) retains barriers \(W=1/4,3/4\), with
\(D_0,N,V\) positive units and bounded derivatives.
Continuous dependence gives
\begin{equation}
 |W(d_q)-W^*(d_q)|\le Cg+C(a_*/d_*)^c\longrightarrow0
 \label{deg:positive-compact-original-buffer}
\end{equation}
uniformly in the independent scales, because \(a_*\le LgkU\).
The original passive-coefficient and raw-buffer margins admit
these terminal values and the regular continuation to the positive port.

\smallskip
\noindent\emph{The initial action factor and the full terminal variations.}
Let \(P_j=\mathcal D_+^jr|_\eta\), \(j=1,2\), and
\(l_j=\mathcal D_+^j\eta_l\). From the original initial height,
\[
 l_1=H(u),\qquad l_2=\mathcal D_+H(u),\qquad
 |l_1|\le C,\quad |l_2|\le Cp,\qquad
 \mathcal D_+r_l=h,\quad |\mathcal D_+h|\le Cp.
\]
Both moving initial rows are
\begin{equation}
 P_{1,l}=h-V_ll_1\asymp p>0,\qquad
 P_{2,l}=\mathcal D_+(P_{1,l})-V_{r,l}P_{1,l}l_1=O(p),
 \qquad
 V_l=\frac{gp(1+\ell_0Q_0(a))}{N_l}.
 \label{deg:positive-compact-inverse-initial-rows}
\end{equation}
The derivative of \(P_{1,l}\) here includes both moving initial
coordinates. The inverse equations \eqref{deg:scalar-inverse-endpoints} give
\((P_2/P_1)_\eta=V_{rr}P_1\).
On the weak box, direct differentiation with \(W_r=-2W\) gives
\[
 V_r\le Cg,\qquad |V_{rr}|\le C(g+W),\qquad
 0<P_1(s)\le Cp e^{Cgs}.
\]
The integral controlling \(P_2/P_1\) is at most
\[
 CgpT e^{CgT}+Cp e^{CgT}\int W\,d\eta\le C
\]
by \eqref{deg:positive-compact-polynomial-absorption} and
\(\int W\,d\eta\le Cw_0\). Thus \(|P_2|\le CP_1\) there.
On \(a_*\le x\le d_*\), the negative ratio \(V_r/V\) gives
\begin{equation}
 0<P_1(x)\le Cp e^{CgT}(a_*/x)^\alpha,\qquad
 |P_2(x)|\le CP_1(x).
 \label{deg:positive-compact-inverse-decay}
\end{equation}
The second bound follows by integrating \(|V_{rr}|P_1\,d\eta\)
with \(V\) a unit. At \(d_*\) both rows are \(O(a^\alpha)\).
The remaining fixed segment has bounded \(\eta\)-length, since
\[
 \eta_T-\eta_*=2\log(d_q/d_*)+\log(W_T/W_*).
\]
The barriers bound the last logarithm, so the inverse equations
transfer both \(O(a^\alpha)\) rows to the original cut.

At the original fixed terminal \(x=d_q\), both full implicit rows
\eqref{deg:scalar-inverse-endpoints}, with \(D=\mathcal D_+\),
give \(|t_1|+|t_2|\le Ca^\alpha\le Cg\), since \(V_T\) is a unit.
The initial factor \(p\) absorbs the arbitrary \(|\log p|\)
weak dwell. The weak-stage bound \(V_r\le Cg\) suffices, so this
argument also applies in the compact ghost box.

\smallskip
\noindent\emph{The exact moving factor and differentiated density.}
Use \eqref{deg:quadratic-outer-moving-factor}--\eqref{deg:quadratic-outer-density}
with \(s=+1,x_0=u,p=P(u)/u^2\). These are exact identities of
this same field; only their present bounded-domain estimates are used.
Write \(\mathscr B=\mathscr B_+\), \(\mathscr J=J_+\),
\(\nu=6g/(1+2g)\), \(\lambda=2\nu\). Direct rational differentiation gives
\(|\mathcal D_+^j\log(-\mathscr B)|\le Cgp\), \(j=1,2\).
At fixed \(\eta\), \(W_r=-2W\), and direct differentiation gives
\[
 |\partial_r^j\mathscr J|\le C(x+kg^2/x),\quad 0\le j\le2,
 \qquad |\mathscr J_\eta|\le CxW.
\]
Let
\[
 I(u)=\int_{\eta_l(u)}^{\eta_T(u)}
            \mathscr J(r(u,\eta),\eta)\,d\eta.
\]
Its two complete action rows are \eqref{deg:scalar-density-endpoints},
with \(D=\mathcal D_+\), \(Dr_l=h\), and the fixed terminal \(r_T\).
On the weak box, \(x+kg^2/x\le Cg\).
The integral bounds just established give
\[
 CgpT e^{CgT}\le CgL_g
\]
for every integral in \eqref{deg:scalar-density-endpoints};
here \(P_1^2\le CP_1\) and \(|P_2|\le CP_1\).
On \(a_*\le x\le d_*\), the common bound is
\[
 Cp e^{CgT}\int_{a_*}^{d_*}
 (x+kg^2/x)(a_*/x)^\alpha\,d\log x
 \le Cp e^{CgT}(a_*+kg^2/a_*)\le Cg .
\]
This uses \(\alpha>1\), \(a_*\asymp a\), and
\(kg^2/a=g/u\le Cg\).
On the remaining fixed strong segment the integrals are bounded
by \(Ca^\alpha\).
The initial boundary terms are \(O(g)\); the full terminal terms
are \(O(a^\alpha)\) by
\eqref{deg:scalar-inverse-endpoints}.
Hence
\begin{equation}
 |\mathcal D_+I|+|\mathcal D_+^2I|\le CgL_g.
 \label{deg:positive-compact-density-bound}
\end{equation}
At artificial splits the full Leibniz terms cancel in this single integral.

\smallskip
\noindent\emph{The complete slope on the original observer image.}
Use~\eqref{deg:quadratic-outer-exact-slope} with \(s=+1,x_0=u\),
including the regular continuation in \(\Phi\) on the admitted
image~\eqref{deg:positive-compact-original-buffer}. Normalize its
coefficient-only reference slope as in that formula; the original
ordinary \(C^3\) logarithmic-slope bound still holds. The initial rows are
\[
 \mathcal D_+\log p=Ep/B=O(1),\qquad
 \mathcal D_+^2\log p=O(p),\qquad
 \mathcal D_+\log a=h,\quad \mathcal D_+^2\log a=\mathcal D_+h=O(1).
\]
Thus the powers \(p^{-\nu}\) and \(a^{-\lambda}\) contribute
\(O(g)\); the rows of \(\log(1+\ell_0Q_0(a))\) and
\(\log(-\mathscr B)\) are \(O(gp)\).
The full terminal formulas~\eqref{deg:scalar-inverse-endpoints} and
the ordinary \(C^3\) observer bound give \(O(a^\alpha)\) for both
logarithmic rows of the terminal factors. Combining these estimates
with~\eqref{deg:positive-compact-density-bound} proves
\eqref{deg:positive-compact-outer-jets}, uniformly as \(p\to0\).
The exact slope retains \(d_\theta=c_\theta(gk)^2\ell_0^{-1-\nu}g^{-\nu}\);
frozen affine calibration preserves its logarithmic rows.
\end{proof}

\subsection{Separate inward jets in the joint quadratic corner}

The derivative of the curvature error involves second derivatives of
the relative matching slope. To control these when both lips grow,
we construct three action derivatives of each inward remainder on its
own domain. Their quotient then has the required mixed derivatives on
the product, before equality of the two actual central heights is imposed.

\begin{lemma}\label{deg:lem-joint-inward-product-direct}
Use the original field \eqref{deg:quadratic-field}, its moving lips
\(Y=\ell_0P\), the original central section \(\xi=0\), and a fixed
compact passive \(A,C\) neighborhood. Fix the original sufficiently
small \(\delta_q>0\), \(\pi\in[\pi_\flat,\Pi]\), and a sufficiently
large fixed \(M\). For independently small \(g,k>0\), restrict the
positive side to strict root-free fibres \(\Delta<0\). On each separate
original admissible lip domain with
\[
 u,r\ge2M,\qquad ku,kr\le\delta_q,
\]
the positive inward passage and the inverse negative passage reach
the same original central section at actual positive heights
\(Z_+(u)\), \(Z_-(r)\). Use \eqref{deg:quadratic-inward-family}--\eqref{deg:quadratic-inward-memory-common}
with \(x_+=u,x_-=r\), \(\theta_s=x_s^{-1}+kx_s\).
The separate central remainders satisfy
\[
 \sum_{j=1}^3|\mathcal D_s^j\mathcal R_s|\le Cg\theta_s.
\]
Their values retain any original coefficient-only logarithmic constants;
no small zeroth remainder row is asserted.

On the product of these separate actual domains define
\[
 N_\pm=-g\mathcal D_\pm\log Z_\pm
       =1-gH_\pm-g\mathcal D_\pm\mathcal R_\pm,\qquad
 H_\pm=P_\pm'/(xB_\pm),\qquad n=N_+/N_-,\quad\varrho=n-1 .
\]
Then \(N_\pm\ge1/2\) and
\begin{equation}\label{deg:joint-inward-product-direct-jets}
 \sum_{i+j\le2}|\mathcal D_+^i\mathcal D_-^j\varrho|
                 \le Cg(\theta_++\theta_-).
\end{equation}
On the original actual subset \(Z_+(u)=Z_-(r)\), this \(n\) is exactly
the relative variation in
\[
 R_u=\frac{uB_+(u)\widehat P(R)}{RB_-(r)P_+(u)}\,n,
 \qquad R=kr,\quad \widehat P(R)=k^2P_-(R/k).
\]
The admissible central-image intersection may be empty. All coefficients
are frozen under the phase operators; constants are independent of the
ghost gap and of any comparison between \(g\) and \(k\).
\end{lemma}

\begin{proof}
Choose a fixed small \(0<b<M\), with \(\Pi b<1/4\), so
\(P_s\ge p_b>0\) on \([0,b]\) after the independent \(k\)-shrink.
On the positive strict ghost path and the negative joint interval,
\(P_s>0\); \(B_s,Q_s\) are positive units. Decreasing \(v\)
follows the positive passage or the inverse negative passage,
respectively, in~\eqref{deg:quadratic-inward-scalar-common}.
For \(v\ge b\), put \(t=v^{-1}\), \(V=kv\), \(\vartheta=t+V\), and
\[
 \begin{gathered}
 D_s=\frac{3-s\pi t+t^2+sV}{1+2g+sgV}
                  (-t\partial_t+V\partial_V),\\
 h_s=P_s/(v^2B_s),\quad H_0=\frac6{1+2g},\quad\ell_g=1-Ag+Cg^2.
 \end{gathered}
\]
These variables lie in a fixed compact rectangle and \(h_s\le C_b\),
including at the ghost bottleneck. The exact formulas
\[
 H_s-H_0=\frac{s[-\pi t+(3-gH_0)V]}{B_s},\qquad
 \omega_s=\frac{\zeta_sP_s}{vB_s}
                 =sCg^3V(1+sV)h_s
\]
give, without dividing by \(P_s\),
\begin{equation}\label{deg:joint-inward-weighted-coefficients}
 \begin{aligned}
 |D_s^j(H_s-H_0)|&\le C\vartheta &&(0\le j\le3),\\
 |D_s^j(Q_s-\ell_g)|&\le CgV,\quad
 |D_s^j\omega_s|\le Cg^3V &&(0\le j\le3),\\
 |D_s^j\vartheta|&\le C\vartheta &&(1\le j\le3).
 \end{aligned}
\end{equation}
Use the ratio equation and exact memory of
Lemma~\ref{deg:lem-stable-ratio-ledger}, with
\[
 \tau=(A_s(x_s)-A_s(v))/g,\quad
 T=(A_s(x_s)-A_s(b))/g,\quad W=Y/P_s(v),\quad W(0)=\ell_0,
\]
\[
 f=(Q_s-\omega_s)W/(1+Q_sW),\qquad
 I_b=\int_0^{T(x_s)}f\,d\tau .
\]
At fixed \(\tau\), \(\mathcal D_{x_s}v=P_s(v)/(vB_s(v))\), so
phase rows transport to powers of \(D_s\); clock rows use \(-gD_s\).
They commute, and \(|\vartheta_\tau|\le Cg\vartheta\) gives
\(\vartheta(v(\tau))\le\theta_s e^{Cg\tau}\).
Thus~\eqref{deg:joint-inward-weighted-coefficients} verifies all
phase and mixed coefficient inputs with
\(D=\mathcal D_{x_s}\), \(\kappa=g\), \(\rho=g\theta_s\),
\(\alpha=Cg<c_0/12\). The fixed small lip gives the two negative
margins. The lemma constructs the actual graph to \(b\) and bounds
its fixed-clock rows by \(Cg\theta_s(1+\tau)^Ne^{-c\tau}\).
Each strict fibre has finite action and flight time; their lengths
need no uniform upper bound.

The moving clock has \(\mathcal DT=1/g\), higher rows zero, and
\[
 T\ge c_b\log(x_s/b)/g\ge a_b/g,\qquad
 g^{-N}(1+T)^Ne^{-cT}\le Cg/x_s
\]
for every fixed \(N\), after shrinking \(g\). The first lower bound
absorbs \(x_s\); the second absorbs every fixed power of \(g^{-1}\).
Hence the full terminal formulas~\eqref{deg:stable-ratio-terminal-rows}
give \(\sum_{j=1}^3|\mathcal D^jI_b|\le Cg\theta_s\).
Their exact entrance memory is
\[
 \begin{gathered}
 y_b=\ell_0P_s(x_s)e^{-T+I_b}=P_s(b)W(T)\le Ce^{-cT},\\
 \Lambda_1=\mathcal D\log y_b=H_s(x_s)-1/g+\mathcal DI_b,\\
 \Lambda_j=\mathcal D^j\log y_b=D_s^{j-1}H_s+\mathcal D^jI_b
 \quad(j=2,3).
 \end{gathered}
\]
Thus \(|\Lambda_1|\le C/g\) and
\(|\Lambda_2|+|\Lambda_3|\le C\theta_s\).
On \([0,b]\), preparation~\eqref{deg:positive-large-S-inward-compact-graph}
checks the fixed positive margin, the \(C/g\) height bounds,
\(y_b\le h_b/2\), and small \(y_b/g\).
Lemma~\ref{deg:lem-compact-height-continuation} therefore gives the
actual centre height, with multiplier \(e^{-A_s(b)/g}\), and
\[
 \log Z_s=\log\{\ell_0P_s(x_s)\}-A_s(x_s)/g+I_b+\mathcal N_s(y_b).
\]
Its three composition rows are bounded by \(Cy_b/g^4\le Cg/x_s\).
This proves the weighted remainder estimate on the original centre
and admits its height to the original low-height window. Concatenation
at the same lips retains the original raw images.

\smallskip
\noindent\emph{Separated product and actual matching.}
For the product estimate let \(N_0=1-gH_0\) and \(d_s=N_s-N_0\).
The weighted coefficient and remainder rows give
\[
 \sum_{j=0}^2|\mathcal D_s^jd_s|\le Cg\theta_s,\qquad
 N_s\ge1/2,\qquad
 \varrho=\frac{d_+-d_-}{N_-}.
\]
The common homogeneous term \(N_0\) cancels in this quotient.
The quotient rule through order
two bounds all separated product rows; in particular the genuinely
mixed row is
\[
 \mathcal D_+\mathcal D_-n
 =-\frac{(\mathcal D_+N_+)(\mathcal D_-N_-)}{N_-^2}.
\]
This proves \eqref{deg:joint-inward-product-direct-jets} on the
product before imposing a central-image intersection.
Equality of the actual central heights
then gives
\[
 r_u=\frac{uB_+P_-(r)}{rB_-P_+(u)}\,\frac{N_+}{N_-},
\]
which is precisely the asserted relative variation after \(R=kr\).
Only the original central-image intersection is used; no actual image
or physical carrier is enlarged.
\end{proof}

\subsection{Physical resection at two simple quadratic roots}

Near \(b/a=2\), the endpoint value of the curvature principal part
can vanish at the same time as both action speeds. Fixed normal sections
replace the lip's endpoint powers by pure action residues. This yields
errors that vanish with each independent footpoint distance. The actual
resection below then transfers the resulting curvature orientation back
to the moving-lip description of the same physical map.

\begin{lemma}[Joint simple-root resection on the physical sections]
\label{deg:lem-joint-root-physical-resection}
Use the positive quadratic field and its original raw section buffers.
Restrict the real-root coefficients to a fixed neighborhood of
\(b/a=2\), so that \(b-a\) is a positive unit. All phase derivatives
freeze \(g,k,\pi,A,C\). Write
\[
 P(u)=(\mathfrak c+ku)(u-a)(u-b),\qquad
 \widehat P(R)=(\mathfrak c-R)(R+ka)(R+kb),
\]
\[
 B_+(u)=1+g(2+ku),\qquad B_-(R)=1+g(2-R),
\]
\[
 \mathfrak a_+=\frac{aB_+(a)}{(\mathfrak c+ka)(b-a)},\qquad
 \mathfrak a_-=
 \frac{\mathfrak cB_-(\mathfrak c)}
      {(\mathfrak c+ka)(\mathfrak c+kb)}.
\]
The residues and their reciprocals are bounded. Keep all
coefficient-only constants in the original actions
\(A_+=J_+\), \(A_-=A_-(R)\). Their exact pole subtractions are
\[
 A_+(a-h)=-\mathfrak a_+\log h+j_+(h),\qquad
 A_-(\mathfrak c-t)=-\mathfrak a_-\log t+j_-(t),
\]
where \(j_s-j_s(0)\) have uniform \(C^3\) bounds.

There are fixed normal-height sections in \(y_+=gY\) and
\(y_-=gk^2Y\), and fixed \(h_0,t_0,g_0,k_0>0\), with the
following properties for independently small \(g,k>0\).
The normal-section axis footpoints are \(u=a-h\) and
\(R=\mathfrak c-t\). Each normal section has its own actual outward
map to the original raw port and observer image.
On the actual intersection of the matched graph with
\(0<h<h_0\), \(0<t<t_0\),
\begin{equation}
 \begin{gathered}
 A_+(a-h)-A_-(\mathfrak c-t)=C_{g,k}+E_+(h)+E_-(t),\\
 |(h\partial_h)^jE_+|\le C\varepsilon_g h,\qquad
 |(t\partial_t)^jE_-|\le C\varepsilon_g t,\quad 0\le j\le3.
 \end{gathered}
 \label{deg:joint-root-normal-weighted-matching}
\end{equation}
where \(E_+(0)=E_-(0)=0\) and
\(\varepsilon_g=Cg(1+|\log g|)^4\).
No bound on \(C_{g,k}\) is required.
The quotient \(n\) of the two separated central action derivatives
has
\begin{equation}
 \begin{gathered}
 \sum_{i+j\le2}
 |\mathcal D_+^i\mathcal D_-^j(n-1)|
 \le C\varepsilon_g(h+t),\\
 \sum_{Z\in\{n,\mathcal D_+n,\mathcal D_-n\}}
 (|\partial_hZ|+|\partial_tZ|)\le C\varepsilon_g .
 \end{gathered}
 \label{deg:joint-root-normal-matching-rows}
\end{equation}
where
\[
 \mathcal D_+=\frac{P(u)}{uB_+(u)}\partial_u,\qquad
 \mathcal D_-=\frac{\widehat P(R)}{RB_-(R)}\partial_R.
\]
The normal-section embeddings use labels
\(\beta_+=-u^2/2\), \(\beta_-=-R^2/2\). After separate
coefficient-only positive affine value and slope calibrations,
their complete maps satisfy
\begin{equation}
 \begin{gathered}
 \phi_s=\mathrm{id}+O_{C^3}(\varepsilon_g),\\
 |\mathcal D_+^j\log\phi_+'|\le C\varepsilon_g h,\qquad
 |\mathcal D_-^j\log\phi_-'|\le C\varepsilon_g t,\quad j=1,2.
 \end{gathered}
 \label{deg:joint-root-normal-outer-rows}
\end{equation}
These statements use only the original ordinary \(C^3\)
logarithmic-slope reserve on the terminal observers, including the
divided finite closing on the negative side.

Let \(x=a-u_{\rm lip}\), \(y=\mathfrak c-R_{\rm lip}\) on the
original lips. The actual resection to the normal sections is
increasing and satisfies
\begin{equation}
 \begin{gathered}
 \mathfrak a_+\log h-(\mathfrak a_++g)\log x=C_++O(x+h),\\
 \mathfrak a_-\log t-(\mathfrak a_-+g)\log y=C_-+O(y+t),
 \qquad |C_s|\le C .
 \end{gathered}
 \label{deg:joint-root-exact-lip-resection}
\end{equation}
The remainders have the corresponding finite logarithmic derivative
bounds. The negative constant contains no unbounded pure \(\log k\).
If \(F_s\) denotes this resection in the internal labels, then
\begin{equation}
 S_{\rm fp}=F_-\circ S_{\rm lip}\circ F_+^{-1},\qquad
 \phi_{{\rm lip},s}=\phi_{{\rm fp},s}\circ F_s
 \label{deg:joint-root-physical-composition}
\end{equation}
on the actual common image, up to the retained positive affine
calibrations of the physical coordinates. Thus both constructions
describe the same completed physical map. No affine identification
of the two internal root-distance coordinates is asserted.
\end{lemma}

\begin{proof}
\emph{The normal sections and their exact invariant.}
On the positive side reverse time and use \(y=gY\); on the
negative side multiply time by \(k\) and use \(y=gk^2Y\).
Division by the positive axis normal rate gives, on both root boxes,
\begin{equation}
 v'=gF_s(v)+c_s(v)y,\qquad y'=y\{1+\kappa_s(v)y\},
 \label{deg:joint-root-local-normal-field}
\end{equation}
where
\[
 \begin{array}{lll}
 F_+=P/(vB_+),&
 c_+=Q_+/(vB_+),&
 \kappa_+=Ckg^2(1+kv)/(vB_+),\\
 F_-=\widehat P/(vB_-),&
 c_-=Q_-/(vB_-),&
 \kappa_-=-Cg^2(1-v)/(vB_-),
 \end{array}
\]
and \(Q_-=1-Ag(1-v)+Cg^2(1-v)^2\).
At \(r_+=a\), \(r_-=\mathfrak c\), the axis rates are exactly
\(gF_s'(r_s)=-g/\mathfrak a_s\).
On fixed boxes bounded away from zero, \(gF_s\) has \(C^4\) norm
\(O(g)\), \(c_s=1/v+O_{C^4}(g)\), and
\(\kappa_s=O_{C^4}(g^2)\), uniformly in \(k\).

Apply the backward-integral construction
\eqref{eq:saddle-normal-time}--\eqref{eq:exact-saddle-invariant}
to \(h=gF_s\). The coefficient bounds just established and
\(Y(-T)\le Cy e^{-T/2}\) give differentiated kernels
\(C(1+T)^Ne^{-(1/2-Cg)T}\) through order three.
This supplies \(\pi_s\) and \(Z_s=ye^{-\mathcal B_s}\), with
\(\pi_s'=gF_s(\pi_s)\), \((\log Z_s)'=1\).
The same estimates for the difference from
\(v'=y/v,\ y'=y\) give
\begin{equation}
 \begin{gathered}
 \pi_s(v,y)=\sqrt{v^2-2y}+O_{C^3}(g),\qquad
 \|\mathcal B_s\|_{C^3}\le Cg^2,\\
 \mathcal B_s=O(g^2y),\qquad \partial_v\pi_s\ge\tfrac12.
 \end{gathered}
 \label{deg:joint-root-local-projection}
\end{equation}
Consequently \(\log Z_s-A_s(\pi_s)/g\) is constant along each orbit.

Choose fixed small heights \(y_{0,s}>0\). On \(y=y_{0,s}\) write
\(\pi_s(N_s(U),y_{0,s})=U\). Then
\[
 N_s(U)=\sqrt{U^2+2y_{0,s}}+O_{C^3}(g),\qquad
 Z_{N,s}(h)=y_{0,s}e^{-\mathcal B_s(N_s(r_s-h),y_{0,s})},
\]
\[
 |\partial_h^j\log Z_{N,s}(h)|\le Cg^2,\qquad 1\le j\le3.
\]
These are the fixed normal sections and their actual axis footpoints.
The physical orbit from a sufficiently small lip point reaches its
normal section inside the unchanged box: the normal height grows
at a fixed rate, the projected axis footpoint remains on the inner
side of the root, and its physical displacement is \(O(y_{0,s})\).

\emph{The local amplitudes and the middle first multipliers.}
Fix inner cuts \(\bar v_s=r_s-\delta_s\) in the two boxes.
The invariant gives the exact implicit relation
\[
 \log y_{\rm in}+e_s(y_{\rm in})
 =\log Z_{N,s}(h)-\{A_s(r_s-h)-A_s(\bar v_s)\}/g,
\]
\[
 e_s(y)=-\mathcal B_s(\bar v_s,y)
 -\{A_s(\pi_s(\bar v_s,y))-A_s(\bar v_s)\}/g .
\]
Here \(e_s(0)=0\), and its required height derivatives have
polynomial \(g^{-1}\) bounds, since the inner cut is a fixed
distance from its root. Solving this scalar identity gives
\begin{equation}
 \begin{gathered}
 y_{{\rm in},s}(h)=\alpha_{y,s}h^{\mathfrak a_s/g}
 e^{-[j_s(h)-j_s(0)]/g}u_s(h),\\
 u_s(h)=\frac{Z_{N,s}(h)}{Z_{N,s}(0)}
 e^{-e_s(y_{{\rm in},s}(h))}.
 \end{gathered}
 \label{deg:joint-root-local-pure-residue}
\end{equation}
with
\[
 \alpha_{y,s}=Z_{N,s}(0)
 e^{[A_s(\bar v_s)-j_s(0)]/g},\qquad
 \alpha_{y,s}^{\pm1}\le e^{K/g}.
\]
The coefficient-only constants cancel in
\(A_s(\bar v_s)-j_s(0)\); thus this local bound permits a
global \(\log k\) term in the negative action. Passing to \(Y\)
on the positive cut and \(\widehat Y=k^2Y\) on the negative cut
adds only the coefficient-only factor \(g^{-1}\).

Choose preliminary \(h_0,t_0\) so the positive residue powers dominate these
fixed exponential constants. Three ordinary endpoint derivatives
introduce only polynomial \(g^{-1}\) factors and inverse endpoint
powers. Their nonlinear terms are then bounded by arbitrarily high
fixed endpoint powers after shrinking \(g\). The remaining phase
part of \(\log u_s\) is \(O_{C^3}(g^2)\).
Consequently the first three logarithmic endpoint rows of
\(g\{\log u_s(h)-\log u_s(0)\}\) are bounded by
\(C\varepsilon_gh\).

For the negative growing bridge, take \(R_0=\bar v_-\), the
already fixed inner cut with a margin below the root and away from
zero, and decrease \(k_0\)
so that \(kr_c<R_0\). Retain its exact first multiplier
\[
 \widehat m=\exp\left\{-\frac1g\int_{kr_c}^{R_0}
       \frac{RB_-}{\widehat P}\,dR\right\},\qquad
 \widehat m/k^2\le e^{K/g}k^{c_0/g-2}\le e^{K/g}.
\]
Indeed \(\widehat P\asymp R^2\) on \(Kk\le R\le R_0\), giving
the action lower bound \(c_0\log(1/k)-K\). Choose \(g_0<c_0/4\).
No reciprocal bound for this collapsing multiplier is used.
Put \(W=\widehat Y/\widehat P\) and
\(q=(A_-(R_0)-A_-(R))/g\). Coefficients and both radial cuts are
fixed under the incoming height variation; the flight
\(T=(A_-(R_0)-A_-(kr_c))/g\) may be arbitrarily long.
On \(kr_c\le R\le R_0\), the fixed root margin gives
\(\widehat P\asymp(R+k)^2\) and \(k/R\le1/r_c\). Hence
\[
 Q_--1=O(g),\qquad
 H=\widehat P_R/(RB_-)=O(1),\qquad
 \omega=-Cg^3(1-R)\widehat P/(RB_-)=O(g^3).
\]
The exact ratio equation has \(\kappa=g\). On the unchanged fixed
small buffer \(0\le W\le\ell_0\), these bounds keep
\(1+Q_-W\) a positive unit and both
\(\mathcal F/W,\mathcal F_W\le-c_1<0\), after decreasing \(g_0\).
At the fixed initial cut \(\widehat P(R_0)\) is a positive unit.
The residue formula \eqref{deg:joint-root-local-pure-residue},
with its factor \(g^{-1}\), places
\(w(t)=\widehat Y_0(t)/\widehat P(R_0)\le\ell_0/2\)
after the same independent coefficient shrink. Thus
Lemma~\ref{deg:lem-fixed-route-height-rows} applies on this full buffer.
Retaining the terminal polynomial in its exact cocycle gives
\[
 \widehat Y_T=\widehat Y_0\widehat m
             e^{\mathcal N_T(w)},\qquad
 |\partial_w^j\mathcal N_T|\le C\quad(1\le j\le3),\qquad
 |(w\partial_w)^j\mathcal N_T|\le Cw\quad(0\le j\le3).
\]
The normalized nonlinear unit therefore has bounded ordinary \(C^3\)
rows in the incoming height at the fixed cut, uniformly in \(g,k,T\).
Its three endpoint rows follow by the Euler composition formula of
Lemma~\ref{deg:lem-compact-height-continuation}, with
\(\Lambda_j=(t\partial_t)^j\log w\): the flat \(w\) absorbs the
polynomial \(g^{-1}\) and endpoint powers from the local residue.
These are fixed-route height rows; phase and moving-terminal rows on
other routes remain those of Lemma~\ref{deg:lem-stable-ratio-ledger}.
Conversion to \(Y\) contributes \(k^{-2}\) only to
\(\widehat m/k^2\), with no reciprocal gain estimate.

Fix the prescribed height buffers of the final compact inward
continuations and their coefficient bounds before the final choice
of \(h_0,t_0\). Enlarge \(K_{\rm tot}\) to include the regular
local action parts and the preceding bound for \(\widehat m/k^2\).
For any fixed endpoint power \(n\) and finite polynomial loss
\(g^{-M}\), choose \(t_0<1\) so that
\(\mathfrak a_{\min}\log(1/t_0)>K_{\rm tot}+2c_*\), with
\(c_*>0\). Then decrease \(g_0\); the residue and its three
ordinary endpoint derivatives give
\[
 g^{-M}e^{K_{\rm tot}/g}t^{\mathfrak a_{\min}/g-j}
 \le Ce^{-c_*/g}t^n,\qquad 0<t\le t_0,\quad 0\le j\le3.
\]
In particular the converted input satisfies \(Y_T\le h_b/2\)
and small \(Y_T/g\) in its already prescribed buffer \(h_b\),
uniformly in independently small \(k\). The positive input has
the same bound after the analogous final choice of \(h_0\).
Lemma~\ref{deg:lem-compact-height-continuation} now applies to the
actual final inward segments: their linear propagators contract,
and their nonlinear height bounds are polynomial in \(g^{-1}\).
Its three Euler composition rows and the displayed superpower
bound control the endpoint rows on those original centre windows.
No reciprocal linear multiplier is introduced in this admission.
Equality of the two exact original central heights therefore gives
\eqref{deg:joint-root-normal-weighted-matching}; all coefficient-only
first multipliers remain in \(C_{g,k}\).
Put \(d_+=h\), \(d_-=t\). On the fixed root boxes,
\(\mathcal D_s=-d_s\gamma_s(d_s)\partial_{d_s}\), where
\(\gamma_s\) are positive units with bounded \(C^2\) norms.
The corrected actions \(A_+-E_+\) and \(A_-+E_-\) in
\eqref{deg:joint-root-normal-weighted-matching} have action derivatives
\[
 \mathcal N_+=1-\mathcal D_+E_+,
 \qquad \mathcal N_-=1+\mathcal D_-E_-,
 \qquad n=\mathcal N_+/\mathcal N_- .
\]
The three weighted rows of \(E_s\) give
\[
 |\mathcal D_s^j(\mathcal N_s-1)|
     \le C\varepsilon_gd_s\quad(0\le j\le2),\qquad
 |\partial_{d_s}\mathcal N_s|
  +|\partial_{d_s}(\mathcal D_s\mathcal N_s)|
     \le C\varepsilon_g .
\]
For the ordinary derivatives, divide the second and third action
rows of \(E_s\) by the single factor \(d_s\gamma_s\).
After shrinking the fixed boxes, \(\mathcal N_s\ge1/2\).
The separated quotient calculation in the proof of
Lemma~\ref{deg:lem-joint-inward-product-direct} now applies with
errors \(\varepsilon_gh\) and \(\varepsilon_gt\); its mixed row is
\(O(\varepsilon_g^2ht)\). Differentiating the quotient and its first
action rows with the displayed ordinary bounds proves
\eqref{deg:joint-root-normal-matching-rows} on the product, before
restriction to the actual match and without comparing \(h\) and \(t\).

\emph{The actual fixed-normal exterior embeddings.}
Put \(\sigma_+=gk\), \(\sigma_-=g\). At a normal section,
\[
 x_l=\sigma_sN_s(U),\qquad z_l=\sigma_s^2y_{0,s}.
\]
The ratio \(W_l=y_{0,s}/N_s(U)^2\) is a positive unit, even
when the normal chord lies beyond its axis root.
Use the signed coefficients of~\eqref{deg:quadratic-outer-ratio},
writing \(N_s^{\rm raw}\) for its numerator and \(D_s=m_s+Q_sW\).
As in the bounded positive outward construction,
\((D_s)_\eta=Q_sW>0\) at \(D_s=0\), so the inverse-height
graph retains its positive denominator.
Choose a fixed auxiliary \(d_*<d_q\) small compared with the lower
bound for \(W_l\). On \([x_l,d_*]\), the singular terms in
\(m_s-sx\) are \(O(g)\). The same invariant-strip argument, starting
now at \(x_l\) with no weak stage, gives
\[
 w_*/2\le W\le w^*,\qquad D_s\ge cW,\qquad
 V_r/V=-2+O((g+x)/w_*)\le-3/2,
\]
\[
 |W(d_*)-W_s^*(d_*)|\le Cg+C(x_l/d_*)^c .
\]
The fixed parabolic continuation therefore admits the graph to its
original raw buffer and observer.

Here \(\partial_U\eta_l=\partial_U^2\eta_l=0\), and the inverse
initial rows are bounded with positive first row. Applying
\eqref{deg:scalar-inverse-endpoints} as above gives
\[
 0<P_1\le C(x_l/x)^{3/2},\quad |P_2|\le CP_1,\qquad
 |\partial_U\eta_T|+|\partial_U^2\eta_T|\le C\sigma_s^{3/2}.
\]
The exact first-variation cocycle and density split, with
\(\nu=6g/(1+2g)\), \(\lambda=2\nu\), give
\[
 J_z=\frac{W_T/D_T}{W_l/D_l}(z_T/z_l)^\nu e^{I_s}.
\]
The density's first two fixed-height radial derivatives are bounded
by \(C(x+kg^2/x)\). Thus~\eqref{deg:scalar-density-endpoints},
with \(D=\partial_U\) and both lower clock rows zero, bounds its two
complete rows by
\[
 C(\sigma_s+kg^2/\sigma_s)+C\sigma_s^{3/2}\le Cg.
\]

Since \(z_l\) is phase-constant and \(\beta=-U^2/2\), the full
moving first variation is
\[
 \Gamma_s'(\beta)=
 \sigma_s^2\frac{N_s(U)N_s'(U)}U
 \frac{W_lN_l^{\rm raw}}{D_l}J_z\Phi'(z_T).
\]
Removing only the coefficient-only observer gain gives the exact
form
\begin{equation}
 \Gamma_s'(\beta)=
 c_s^\ast\sigma_s^{2-\lambda}y_{0,s}^{-\nu}
 \frac{N_s(U)N_s'(U)}U N_l^{\rm raw}
 \frac{W_T}{D_T}z_T^\nu\widehat\Phi'(z_T)e^{I_s}.
 \label{deg:joint-root-normal-exact-physical-slope}
\end{equation}
The initial factor equals \(1+O_{C^2}(g)\), because
\eqref{deg:joint-root-local-projection} gives
\(N_s(U)N_s'(U)/U=1+O_{C^2}(g)\).
The terminal factors have two logarithmic rows
\(O(\sigma_s^{3/2})\), using just the original ordinary \(C^3\)
logarithmic-slope observer bound. Thus both ordinary logarithmic
rows of the slope are \(O(g)\).
Subtracting its exact value and calibrating its exact positive
slope at a fixed reference point gives
\(\phi_s=\mathrm{id}+O_{C^3}(g)\).
The action fields contain the factors \(h,t\), proving
\eqref{deg:joint-root-normal-outer-rows}.
The two physical gains
\(\sigma_+^{2-\lambda}\), \(\sigma_-^{2-\lambda}\) are retained
coefficient-only factors in
\eqref{deg:joint-root-normal-exact-physical-slope}; their ratio
and their reciprocals are not estimated.

\emph{Resection of the lip and the same physical curvature.}
On either lip the local normal coordinate is
\[
 y_{{\rm lip},+}=g\ell_0P(a-x),\qquad
 y_{{\rm lip},-}=g\ell_0\widehat P(\mathfrak c-y).
\]
It is \(g\) times a positive unit times the root distance.
The projected lip distance is
\(\widehat x=x(1+O(g))\), respectively
\(\widehat y=y(1+O(g))\).
Applying the exact invariant between the lip and its normal-section
intersection, and subtracting the two pole expansions, gives
\eqref{deg:joint-root-exact-lip-resection}.
The coefficient terms are \(g\log g\), logarithms of local
positive units, and the bounded projection correction
\(\mathfrak a_s\log(\widehat x/x)\).
The negative comparison is in \(gk^2Y\), so it has no pure
\(\log k\). The common growing-bridge multiplier cancels from
this local comparison.
The finite logarithmic derivative estimates follow by differentiating
this invariant on the positive interior of the lip.
In particular
\[
 h=A_+x^{1+g/\mathfrak a_+}e^{O(x+h)},\qquad
 t=A_-y^{1+g/\mathfrak a_-}e^{O(y+t)},\qquad A_s^{\pm1}\le C .
\]
The increasing maps give uniform inclusions between the nested
endpoint neighborhoods;
the lip exponents are \(\mathfrak a_s/g+1\), whereas the normal
residue quotient has no added \(g\).

Both resection coordinates belong to the same orbit, so
\eqref{deg:joint-root-physical-composition} identifies their completed
maps after the retained positive affine physical calibrations.
The normal physical input is increasing in \(h\), since
\(d\{-(a-h)^2/2\}/dh=a-h>0\); the lip physical input is
increasing in \(x\), and the resection is increasing.
Consequently the same completed physical map has the same curvature
zeros and crossing directions in the two descriptions.
\end{proof}

\subsection{Paired first variation on positive quadratic fibres}

Above a comparison orbit the complete derivative is bounded through
the ratio of two variations. Their separate core factors need not be
units, and both exterior gains contain a small common power of \(gk\).
The next lemma retains these factors until their cancellation in the
physical derivative.

\begin{lemma}[Paired first variation above a quadratic comparison]
\label{deg:lem-positive-quadratic-paired-variation}
Fix the original raw sections, nested buffers, ordinary \(C^3\)
regular observers and the divided finite closing of
Lemma~\ref{lem:closing}.
Use the polynomial and physical coordinates of
\eqref{deg:quadratic-field}, and write \(a=d-\sqrt\Delta\) when
\(\Delta\ge0\). The passive compact
\((A,C)\)-set lies inside their prescribed common coefficient
neighbourhood, with the original buffer and observer margins.
Fix a finite coefficient bound \(\Pi\). In the original domain
construction choose \(X\) and the retained ghost width jointly for
the prescribed finite comparison action threshold, and then fix
\(\pi_\flat\). For \(\pi\in[\pi_\flat,\Pi]\), there are fixed
\(\eta_{\max},C_*>0\) and independent upper bounds for \(g,k\)
with the following property on the retained real, fold and ghost fibres.
Use coefficient fibres with the actual \(K_{\rm fin}>0\) and the
original \(0<m\le m_0\) finite-closing admission.
Choose the comparison lip at
\[
 u_* = a-X/2\quad(\Delta\ge0),\qquad
 u_* = d-X/2\quad(\Delta<0),
 \qquad \eta_{\rm comp}=gZ_+(u_*).
\]
The comparison lies strictly inside the enlarged low lip interval
beginning at \(a-X\), respectively \(d-X\). The actual completed
transition in the original affine input satisfies
\[
 0<\Pi'\le C_*
 \qquad\text{throughout}\qquad
 \eta_{\rm comp}\le\eta_0=gY(0)\le\eta_{\max}.
\]
This entire interval is admitted to the original section buffers.
The constants are independent of the root gap and of any relation
between \(g\) and \(k\). After the common choice \(mC_*<1/2\),
\(F'=1-m\Pi'>1/2\) on it.
\end{lemma}

\begin{proof}
Write
\[
 \begin{gathered}
 \eta=gY,\qquad \epsilon_0=gk,\qquad D=gP+Q\eta,\\
 Q=1-Ag(1+k\xi)+Cg^2(1+k\xi)^2,
 \qquad B=1+g(2+k\xi).
 \end{gathered}
\]
The exact field is
\[
 \dot\xi=-D,\qquad
 \dot\eta=-\eta\{\xi B+H_1\eta\},
 \qquad H_1=Ckg^2(1+k\xi),
 \qquad z=\epsilon_0^2\eta.
\]
By \eqref{ann:source-field}, this is the source field with coefficient quartet
\((-g,g\pi,g,\epsilon_0)\); \(\epsilon_0\) is held fixed under
every phase derivative.

\smallskip
\noindent\emph{The actual comparison and the common cuts.}
Use the jointly chosen fixed \(X>0\). The abscissa \(u_*\) ranges over a
fixed positive compact interval, and there is a fixed \(\epsilon>0\)
such that
\[
 P\ge p_*>0\qquad\text{on }[0,u_*+\epsilon].
\]
The factorization gives this margin. The comparison action choice
of Proposition~\ref{deg:prop-positive-quadratic} first fixes \(X\)
by \(J_+(a-X)\ge c\log(\delta_c/X)\), then the ghost width;
the fold has the stronger inverse-distance divergence.

Corollaries~\ref{deg:cor-positive-compact-suppliers} and
\ref{deg:cor-negative-inward-common-centre} give the original inward
cocycles, without changing the centre,
\[
 \log Z_+=\log(\ell_0P)-J_+/g+\mathcal R_+,
 \qquad
 \log Z_-=\log(\ell_0P_-)-A_-/g+\mathcal R_-.
\]
Their zeroth rows on the bounded comparison and the entire strict
negative side are bounded. Indeed their ratio equations give
\(W\le\ell_0e^{-c\tau}\), hence
\(\left|\int (Q-\omega)W/(1+QW)\,d\tau\right|\le C\ell_0\),
and the original compact continuation adds \(O(Y_b/g)\).
The original action constants remain, \(N_-\ge1/2\), and
\(Z_-\downarrow0\) at its simple root. Choose fixed \(r_{\max}>r_c\)
with \(A_-(r_{\max})\) above the bounded comparison action, then shrink
\(k\). Monotonicity and the zeroth cocycles give the unique actual match
\[
 r_*\in[r_c,r_{\max}],\qquad Z_+(u_*)=Z_-(r_*),
 \qquad \eta_{\rm comp}=gZ_+(u_*).
\]

On either comparison lip, \(z/x^2=\ell_0gp\), with \(p\) a
positive unit: on the negative lip \(r_*\) is bounded. Use the
original inverse logarithmic-height graph until it reaches a fixed
small strip height \(w_0\). Its bounded-comparison weak length is
\(O(1+|\log g|)\). The positive weak displacement estimate is
\(O(g(1+|\log g|))+O(w_0)\) in \(\log x\).
Choose \(w_0\) first so small that this strip is entered before
\(\xi=u_*+\epsilon/2\), and then decrease \(g\).
This uses a fixed positive margin of \(P\), and is uniform through
the fold. The negative weak stage reaches its fixed strip by
\(|\xi|\le L_1r_*\), with \(L_1\) fixed.

The outward constructions of
Corollaries~\ref{deg:cor-positive-compact-suppliers} and
\ref{deg:cor-negative-absolute-outer} continue both comparison
traces. Choose one fixed \(L\) larger than those entrance bounds
and the bounded positive factors. On the two coefficient-dependent
but phase-fixed cuts \(\xi=\pm L\), the comparison has
\[
 c_L\le\eta(\pm L)\le C_L,
 \qquad D(\pm L)\asymp\eta(\pm L).
\]
Increasing \(L\), the strip contraction places both values of
\(\eta/L^2\) inside the same small parabolic band. The original
outward continuation, including its weak-stage initial factors,
lands strictly inside the original raw buffers. No raw section or
observer window has been replaced.

\smallskip
\noindent\emph{One upper reference and all intervening inputs.}
Fix \(\eta_{\max}>0\). At \(g=k=0\) the core graph with centre
height \(\eta_{\max}\) is
\(\eta=\eta_{\max}+\xi^2/2\), with strictly positive height and
speed on every fixed core. Ordinary dependence on \([-L,L]\)
constructs its actual upper reference for small independent \(g,k\).
Choose \(L\) large enough that its two values of \(\eta/L^2\)
belong to the same parabolic band as the comparison.

The inverse-height first-exit calculation continues this reference
to an auxiliary raw cut \(d_*<d_q\), where strip contraction gives
\[
 |W(d_*)-W^*(d_*)|\le Cg+C(\epsilon_0L/d_*)^c\longrightarrow0.
\]
The regular segment to the original port has the comparison's limit
\(W^*(d_q)\); the prescribed passive and raw-buffer margins admit it.

Scalar uniqueness orders the graphs between the comparison and
this upper reference. They stay between the two finite positive
core graphs and between their admitted outward tails. Thus every
\(\eta_0\in[\eta_{\rm comp},\eta_{\max}]\) has an actual through
graph on the entire word and remains in the original convex raw
and observer windows. Its central raw height is
\(z(\gamma)=\epsilon_0^2\eta_0\); after the common coefficient
shrink the original low-height and speed restrictions are satisfied.
The whole interval, including nonroot inputs, is consequently used.

\smallskip
\noindent\emph{The paired core variation.}
For
\[
 f=\frac{d\eta}{d\xi}
   =\frac{\eta(\xi B+H_1\eta)}D
\]
use \eqref{deg:scalar-raw-cocycle}, with residual
\([gP'+(Q'+H_1)\eta]/D\).
On \([-L,u_*+\epsilon]\), \(P\) has a fixed positive lower
bound after decreasing \(k\); on the remaining positive core,
the comparison has already entered its invariant strip. Hence
\[
 D_{\rm comp}\ge c g\quad\text{on }[-L,L],\qquad
 \eta/D\le C
\]
on every graph above it. In the part where \(P\) could be negative,
\(\eta\ge c>0\), so \(D\ge c\eta\) after decreasing \(g\).
Since \(P'\) is bounded and
\(Q'=O(gk)\), \(H_1=O(kg^2)\), the remaining integral in the
identity is bounded in absolute value on the full core. The endpoint
quotients \((\eta/D)(\pm L)\) are positive units.

Set
\[
 A_\pm=\frac{\partial\eta(\pm L;\eta_0)}{\partial\eta_0}>0.
\]
Integration gives
\begin{equation}
 \frac{A_-}{A_+}
 =\frac{(\eta/D)(-L)}{(\eta/D)(L)}
   \exp\left\{\int_L^{-L}
          \frac{gP'+(Q'+H_1)\eta}D\,d\xi\right\}
 \in[c,C].
 \label{deg:quadratic-paired-core-ratio}
\end{equation}
This estimate controls the ratio of the two core variations;
separate unit bounds for \(A_\pm\) are not required.
For comparison with a moving lip phase \(u\), the exact incoming
factor is
\[
 \frac{d\eta_0}{du}
  =\eta_0\left\{\frac{P'}P-\frac{uB_+}{gP}
                     +\mathcal R_+'\right\}
  =-\frac{uB_+\eta_0}{gP}N_+.
\]
The core-end derivatives are \(A_\pm\,d\eta_0/du\), so this
complete moving-source factor cancels. The common cuts are phase-fixed.

\smallskip
\noindent\emph{Both original tails and the physical derivative.}
In the source coordinates \((\xi,\eta)\), with quartet
\((-g,g\pi,g,\epsilon_0)\), the displayed parabolic rational
calculation \eqref{ann:tail-first-derivative} gives, on either
admitted tail,
\[
 F_{\sigma,\eta}-\lambda_g/r
   =O(\epsilon_0+g\pi r^{-2}+r^{-3+\nu}),
 \qquad \lambda_g=\frac{12g}{1+2g},\qquad 0<\nu<1.
\]
Its remainder has bounded integral from \(L\) to
\(d_\sigma/\epsilon_0\). The parabolic bands just constructed
justify this application even though the centre height may be
exponentially small. Consequently the
two actual tail first variations are
\[
 J_\sigma=(d_\sigma/(\epsilon_0L))^{\lambda_g}U_\sigma,
 \qquad U_\sigma,U_\sigma^{-1}\le C.
\]
At the original raw ports the complete core-plus-tail variations are
exactly
\begin{equation}
 \frac{dz_\sigma}{d\eta_0}
   =\epsilon_0^{2-\lambda_g}(d_\sigma/L)^{\lambda_g}
       A_\sigma U_\sigma.
 \label{deg:quadratic-paired-raw-slopes}
\end{equation}
All terminal factors lie on the admitted original images.
The original bridges and divided closing, with \(K_{\rm fin}>0\),
\(m\le m_0\), have positive derivative units by Lemma~\ref{lem:closing}.

The common \(\epsilon_0^{2-\lambda_g}=(gk)^{2-\lambda_g}\)
factor cancels in the ratio of the two raw variations. The ratio
\((d_-/d_+)^{\lambda_g}\), the two \(U\)-factors, and the
paired core ratio are bounded units. Composing the original regular
bridges and divided closing therefore proves \(0<\Pi'\le C_*\)
in the same original affine input on the entire stated interval.
Choose \(mC_*<1/2\) only after these constants and buffers have
been fixed. This gives the stated strict sign of \(F'\).
\end{proof}
